\documentclass[UTF8,11pt,a4paper]{ctexart}
\usepackage{amsmath,amssymb,geometry,enumitem,hyperref,array,longtable,booktabs}
\geometry{margin=2.1cm}\pagestyle{plain}\setlist{itemsep=2pt,topsep=3pt}
\hypersetup{colorlinks=true,linkcolor=blue}
\setlength{\emergencystretch}{2em}
\setlength{\abovedisplayskip}{5pt}
\setlength{\belowdisplayskip}{5pt}
\setlength{\abovedisplayshortskip}{2pt}
\setlength{\belowdisplayshortskip}{2pt}
\newcommand{\EN}[1]{\par\noindent \textbf{English.} #1\par}
\newcommand{\FORMULA}[2]{\par\noindent\textbf{公式名称 / Formula name.} #1 / #2\par}
\newcommand{\READ}[1]{\par\noindent \textbf{公式英文读法 / How to read it.} #1\par}
\newcommand{\EX}[2]{\par\noindent \textbf{例子 / Example.} #1\par\EN{#2}}
\newcommand{\ASK}[4]{\par\noindent \textbf{可能问题 / Possible question.} #1\par\EN{#2}\noindent \textbf{回答要点 / Key answer.} #3\par\EN{#4}\medskip}
\newcommand{\APS}[2]{\par\noindent\textbf{中文。} #1\par\noindent\textbf{English.} #2\par}
\newcommand{\DEF}[2]{\par\noindent\textbf{定义 (Definition).} #1\par\EN{#2}}
\newcommand{\THM}[2]{\par\noindent\textbf{定理 (Theorem).} #1\par\EN{#2}}
\newcommand{\COR}[2]{\par\noindent\textbf{推论 (Corollary).} #1\par\EN{#2}}
\newcommand{\PRF}[2]{\par\noindent\textbf{证明思路 (Proof idea).} #1\par\EN{#2}}
\newcommand{\REM}[2]{\par\noindent\textbf{注 (Remark).} #1\par\EN{#2}}
\title{离散数学 APS 面试准备\\[4pt]\large 第 1、3、6、7、8 章：数理逻辑、二元关系、代数、格与布尔代数、图论}
\author{Discrete Mathematics: Bilingual Review}
\date{}
\begin{document}\maketitle
\section*{课程概述 | Course overview}
\APS{本讲义对应方世昌《离散数学》（第三版，西安电子科技大学出版社，2009）第 1、3、6、7、8 章，以\textbf{定义—定理—证明思路}为主线：先给严格定义，再给完整假设与结论的定理，最后用一到两句话交代证明的关键。五章的次序是一步步抽象：第 1 章用命题与谓词把“推理”形式化，并证明真值表、范式、全功能集、归结等判定工具；第 3 章把“比较”与“相同”抽象为二元关系，证明闭包的构造、Warshall 算法以及等价关系与划分的对应；第 6 章抽出载体与运算，用同态、同余、商代数比较代数系统，并得到 Lagrange 定理与群同态基本定理；第 7 章统一偏序格与代数格，证明两种定义的等价，并给出布尔代数与 Stone 表示；第 8 章用顶点和边描述连接关系，给出握手定理、Euler 公式、Hall 定理、四色定理与树的等价刻画。例子每章只保留一两个，用来核验定义与定理，不再重复算例。}{This review follows Chapters 1, 3, 6, 7 and 8 of Fang Shichang's \emph{Discrete Mathematics} (3rd ed., Xidian University Press, 2009), organised as \textbf{definition--theorem--proof idea}: first a strict definition, then a theorem with complete hypotheses and conclusion, and finally a one- or two-sentence account of the key step. The five chapters abstract step by step. Chapter 1 formalizes reasoning with propositions and predicates and justifies truth tables, normal forms, functional completeness and resolution. Chapter 3 abstracts comparison and sameness into binary relations and proves closure constructions, Warshall's algorithm and the correspondence between equivalence relations and partitions. Chapter 6 extracts carriers and operations, compares algebraic systems through homomorphisms, congruences and quotient algebras, and obtains Lagrange's theorem and the fundamental theorem of group homomorphisms. Chapter 7 unifies the order-theoretic and the algebraic lattice, proves the equivalence of the two definitions, and gives Boolean algebras with Stone's representation. Chapter 8 describes connections by vertices and edges and proves the handshake theorem, Euler's formula, Hall's theorem, the four-colour theorem and the equivalent characterizations of trees. Each chapter keeps only one or two examples to verify definitions and theorems, with no repeated computations.}
\begin{center}\small
\begin{tabular}{@{}>{\raggedright\arraybackslash}p{2.0cm}>{\raggedright\arraybackslash}p{5.2cm}>{\raggedright\arraybackslash}p{7.1cm}@{}}
\toprule
\textbf{章 / Chapter} & \textbf{核心 / Core} & \textbf{面试考点 / Interview focus} \\
\midrule
1 数理逻辑 & 命题演算、范式、推理、谓词与量词 / propositional calculus, normal forms, inference, predicates and quantifiers & 真值表与主范式；全功能集；直接、附加前提、反证、归结；量词否定与 $\forall\exists$ 次序 / truth tables and principal normal forms; functionally complete sets; direct, CP and indirect proofs; resolution; quantifier negation and the order of $\forall\exists$ \\
3 二元关系 & 性质、合成、闭包、偏序、等价关系 / properties, composition, closures, partial orders, equivalence & 五条性质的判定；Warshall；哈斯图与确界；等价类、商集、划分 / testing the five properties; Warshall; Hasse diagrams and bounds; classes, quotient sets and partitions \\
6 代数 & 载体与运算、子代数、同态、同余、商代数、群、环域 / carrier and operations, subalgebras, homomorphisms, congruences, quotients, groups, rings and fields & 群公理与子群判定；循环群与陪集；Lagrange；同态核；$\mathbb Z_n$ 何时为域 / group axioms and subgroup tests; cyclic groups and cosets; Lagrange; kernels; when $\mathbb Z_n$ is a field \\
7 格与布尔代数 & 两种格定义、分配格、有补格、布尔代数 / two definitions of lattices, distributive and complemented lattices, Boolean algebras & 格定义的等价性；$M_3,N_5$；补元唯一性；Stone 表示；开关函数 / equivalence of the two definitions; $M_3,N_5$; uniqueness of complements; Stone representation; switching functions \\
8 图论 & 度与连通、矩阵、匹配、平面性、着色、树 / degrees and connectivity, matrices, matching, planarity, colouring, trees & 握手定理；Euler 公式；Hall 条件；Kruskal/Prim；有向树 / handshake theorem; Euler's formula; Hall's condition; Kruskal/Prim; directed trees \\
\bottomrule
\end{tabular}
\end{center}

\section[数理逻辑（第 1 章）]{数理逻辑（教材第 1 章）| Mathematical logic (Chapter 1)}
\subsection{命题、联结词与真值表 | Propositions, connectives and truth tables}
\DEF{\textbf{命题}是或真或假、但不能既真又假的陈述句；其真值记为 $T$（真，$1$）或 $F$（假，$0$）。不能再分解的命题叫\textbf{原子命题}，用命题变元 $P,Q,R,\dots$ 表示。由命题变元经联结词与括号按规则构成的符号串叫\textbf{命题公式}（合式公式），其归纳定义为：(i) 命题常元 $T,F$ 与命题变元是公式；(ii) 若 $A,B$ 是公式，则 $\neg A$、$(A\wedge B)$、$(A\vee B)$、$(A\to B)$、$(A\leftrightarrow B)$ 是公式；(iii) 只有经 (i)(ii) 得到的才是公式。}{A \textbf{proposition} is a declarative sentence that is either true or false but not both; its truth value is $T$ (true, $1$) or $F$ (false, $0$). A proposition that cannot be decomposed is \textbf{atomic} and is denoted by a propositional variable $P,Q,R,\dots$. A symbol string built from variables and connectives by the rules is a \textbf{propositional formula} (well-formed formula), defined inductively: (i) constants $T,F$ and variables are formulas; (ii) if $A,B$ are formulas, then so are $\neg A$, $(A\wedge B)$, $(A\vee B)$, $(A\to B)$ and $(A\leftrightarrow B)$; (iii) only strings obtained by (i)--(ii) are formulas.}
\APS{五个基本联结词是\textbf{否定} $\neg$、\textbf{合取} $\wedge$、\textbf{析取} $\vee$、\textbf{蕴含} $\to$、\textbf{等值} $\leftrightarrow$，其意义完全由下表确定，与语句内容无关；优先级为 $\neg,\wedge,\vee,\to,\leftrightarrow$，同级左结合，最外层括号可省。}{The five basic connectives are \textbf{negation} $\neg$, \textbf{conjunction} $\wedge$, \textbf{disjunction} $\vee$, \textbf{implication} $\to$ and \textbf{equivalence} $\leftrightarrow$, fixed entirely by the table below, independent of content; precedence is $\neg,\wedge,\vee,\to,\leftrightarrow$, equal connectives associate to the left, and outermost parentheses may be dropped.}
\begin{center}\small
\begin{tabular}{@{}ccccccc@{}}
\toprule
$P$ & $Q$ & $\neg P$ & $P\wedge Q$ & $P\vee Q$ & $P\to Q$ & $P\leftrightarrow Q$ \\
\midrule
1 & 1 & 0 & 1 & 1 & 1 & 1 \\
1 & 0 & 0 & 0 & 1 & 0 & 0 \\
0 & 1 & 1 & 0 & 1 & 1 & 0 \\
0 & 0 & 1 & 0 & 0 & 1 & 1 \\
\bottomrule
\end{tabular}
\end{center}
\REM{$P\to Q$ 只排除“前件真而后件假”这一种情况，因此当前件为假时规定为真。这样 $P\to Q\Leftrightarrow\neg P\vee Q$ 才成立，假言推理与拒取式才在所有赋值下可靠。固定译法：“$P$ 仅当 $Q$”“$Q$ 每当 $P$”“$P$ 是 $Q$ 的充分条件”都译为 $P\to Q$；“$P$ 是 $Q$ 的必要条件”译为 $Q\to P$；$P\to Q$ 的逆反命题 $\neg Q\to\neg P$ 与原命题等值。}{$P\to Q$ excludes only the case ``antecedent true and consequent false'', so it is declared true when the antecedent is false; this makes $P\to Q\Leftrightarrow\neg P\vee Q$ and keeps modus ponens and modus tollens sound under all assignments. Fixed translations: ``$P$ only if $Q$'', ``$Q$ whenever $P$'' and ``$P$ is sufficient for $Q$'' are all $P\to Q$; ``$P$ is necessary for $Q$'' is $Q\to P$; the contrapositive $\neg Q\to\neg P$ of $P\to Q$ is equivalent to it.}

\subsection{重言式、可满足性与逻辑等价 | Tautologies, satisfiability and logical equivalence}
\DEF{给公式中每个变元一个真值的映射叫\textbf{赋值}。含 $n$ 个变元的公式恰有 $2^n$ 个赋值。若公式在一切赋值下为真，称为\textbf{重言式}（永真式）；在一切赋值下为假，称为\textbf{矛盾式}（永假式）；至少在一个赋值下为真，称为\textbf{可满足式}。两个公式 $A,B$ 若在一切赋值下真值相同，即 $A\leftrightarrow B$ 为重言式，则称 $A,B$ \textbf{逻辑等价}，记 $A\Leftrightarrow B$。}{A map assigning a truth value to each variable is an \textbf{assignment}; a formula with $n$ variables has exactly $2^n$ assignments. A formula true under every assignment is a \textbf{tautology}, false under every assignment a \textbf{contradiction}, and true under at least one assignment \textbf{satisfiable}. Two formulas $A,B$ are \textbf{logically equivalent}, written $A\Leftrightarrow B$, when they have the same truth value under every assignment, i.e.\ when $A\leftrightarrow B$ is a tautology.}
\THM{\textbf{替换定理}：若 $A\Leftrightarrow B$，则在任意公式中把 $A$ 的若干出现换成 $B$ 所得到的公式与原公式逻辑等价。基本等价式（都是关于任意公式 $P,Q,R$ 的恒等式）如下表，它们是所有等值变换与化简的依据。}{\textbf{Replacement theorem}: if $A\Leftrightarrow B$, then replacing some occurrences of $A$ by $B$ in any formula yields a formula logically equivalent to the original. The basic equivalences (identities holding for all formulas $P,Q,R$) are listed below; they justify all equivalence transformations and simplifications.}
\begin{center}\small
\begin{tabular}{@{}p{8.8cm}p{7.0cm}@{}}
\toprule
\textbf{等价式 / Equivalence} & \textbf{名称 / Name} \\
\midrule
$\neg\neg P\Leftrightarrow P$ & 双重否定 / double negation \\
$P\wedge P\Leftrightarrow P,\quad P\vee P\Leftrightarrow P$ & 幂等律 / idempotence \\
$P\wedge Q\Leftrightarrow Q\wedge P,\quad P\vee Q\Leftrightarrow Q\vee P$ & 交换律 / commutativity \\
$(P\wedge Q)\wedge R\Leftrightarrow P\wedge(Q\wedge R)$，对 $\vee$ 同理 & 结合律 / associativity \\
$P\wedge(Q\vee R)\Leftrightarrow(P\wedge Q)\vee(P\wedge R)$，对偶式同理 & 分配律 / distributivity \\
$\neg(P\wedge Q)\Leftrightarrow\neg P\vee\neg Q,\quad \neg(P\vee Q)\Leftrightarrow\neg P\wedge\neg Q$ & 德摩根律 / De Morgan \\
$P\wedge(P\vee Q)\Leftrightarrow P,\quad P\vee(P\wedge Q)\Leftrightarrow P$ & 吸收律 / absorption \\
$P\to Q\Leftrightarrow\neg P\vee Q\Leftrightarrow\neg Q\to\neg P$ & 蕴含式 / implication and contrapositive \\
$P\leftrightarrow Q\Leftrightarrow(P\to Q)\wedge(Q\to P)$ & 等值式 / equivalence \\
\bottomrule
\end{tabular}
\end{center}
\COR{$P\wedge(P\to Q)\Rightarrow Q$（假言推理）与 $\neg Q\wedge(P\to Q)\Rightarrow\neg P$（拒取式）是两条永真的蕴含式；一般地，$A\Rightarrow B$ 定义为 $A\to B$ 是重言式。}{$P\wedge(P\to Q)\Rightarrow Q$ (modus ponens) and $\neg Q\wedge(P\to Q)\Rightarrow\neg P$ (modus tollens) are valid entailments; in general $A\Rightarrow B$ means $A\to B$ is a tautology.}

\subsection{范式与主范式 | Normal forms and principal normal forms}
\DEF{命题变元或其否定叫\textbf{文字}；有限个文字的合取叫\textbf{基本积}，有限个文字的析取叫\textbf{基本和}。基本积的析取称为\textbf{析取范式}（DNF），基本和的合取称为\textbf{合取范式}（CNF）。若每个基本积（基本和）含全部 $n$ 个变元各一次，则分别称为\textbf{极小项} $m_i$ 与\textbf{极大项} $M_i$：极小项 $m_i$ 恰在第 $i$ 个赋值下为真，极大项 $M_i$ 恰在第 $i$ 个赋值下为假，且 $M_i=\neg m_i$。极小项之析取叫\textbf{主析取范式}（PDNF），极大项之合取叫\textbf{主合取范式}（PCNF）。}{A \textbf{literal} is a variable or its negation; a \textbf{fundamental product} is a conjunction of finitely many literals and a \textbf{fundamental sum} a disjunction of finitely many literals. A disjunction of fundamental products is in \textbf{disjunctive normal form} (DNF), a conjunction of fundamental sums in \textbf{conjunctive normal form} (CNF). If every fundamental product (sum) contains all $n$ variables exactly once, we get a \textbf{minterm} $m_i$ and a \textbf{maxterm} $M_i$: the minterm $m_i$ is true exactly under the $i$-th assignment, the maxterm $M_i$ is false exactly under the $i$-th assignment, and $M_i=\neg m_i$. A disjunction of minterms is the \textbf{principal disjunctive normal form} (PDNF), a conjunction of maxterms the \textbf{principal conjunctive normal form} (PCNF).}
\THM{\textbf{范式存在与唯一性}：每个命题公式 $A$ 都存在与之等价的 DNF 与 CNF。若 $A$ 不是矛盾式，则 $A$ 有唯一的（不计项与文字的次序及重复）主析取范式；若 $A$ 不是重言式，则 $A$ 有唯一的主合取范式。}{Let $A$ be a propositional formula. Then $A$ has an equivalent DNF and an equivalent CNF. If $A$ is not a contradiction, $A$ has a PDNF unique up to the order and repetition of terms and literals; if $A$ is not a tautology, $A$ has a unique PCNF.}
\PRF{存在性：反复用德摩根律与分配律把 $\neg$ 压到变元、把 $\wedge$ 分配到 $\vee$ 上，即得 DNF/CNF。唯一性：$A$ 的真值表是确定的，其值为真的行编号 $I$ 唯一决定 PDNF $=\bigvee_{i\in I}m_i$，值为假的行编号 $J$ 唯一决定 PCNF $=\bigwedge_{j\in J}M_j$；两种主范式还满足 $\bigvee_{i\in I}m_i\Leftrightarrow\neg\bigwedge_{j\notin I}M_j$。}{Existence: repeatedly apply De Morgan's laws to push $\neg$ down to variables and distributivity to move $\wedge$ over $\vee$, obtaining DNF/CNF. Uniqueness: the truth table of $A$ is fixed, so the set $I$ of true rows uniquely determines the PDNF $\bigvee_{i\in I}m_i$ and the set $J$ of false rows uniquely determines the PCNF $\bigwedge_{j\in J}M_j$; moreover $\bigvee_{i\in I}m_i\Leftrightarrow\neg\bigwedge_{j\notin I}M_j$.}

\subsection{联结词的完备性 | Functional completeness}
\DEF{设 $S$ 是若干联结词组成的集合。若任何命题公式都逻辑等价于一个只含 $S$ 中联结词的公式，则称 $S$ 为\textbf{全功能集}（完备集）。除五个基本联结词外，常用\textbf{异或} $P\oplus Q\Leftrightarrow\neg(P\leftrightarrow Q)$、\textbf{与非} $P\uparrow Q\Leftrightarrow\neg(P\wedge Q)$、\textbf{或非} $P\downarrow Q\Leftrightarrow\neg(P\vee Q)$。}{Let $S$ be a set of connectives. If every propositional formula is logically equivalent to a formula using only connectives from $S$, then $S$ is \textbf{functionally complete} (adequate). Besides the five basic connectives one uses \textbf{exclusive or} $P\oplus Q\Leftrightarrow\neg(P\leftrightarrow Q)$, \textbf{NAND} $P\uparrow Q\Leftrightarrow\neg(P\wedge Q)$ and \textbf{NOR} $P\downarrow Q\Leftrightarrow\neg(P\vee Q)$.}
\THM{\textbf{完备性}：$\{\neg,\wedge,\vee\}$ 是全功能集；由德摩根律，$\{\neg,\wedge\}$、$\{\neg,\vee\}$ 与 $\{\neg,\to\}$ 也是全功能集；单元素集 $\{\uparrow\}$ 与 $\{\downarrow\}$ 各自都是全功能集。反之，$\{\wedge,\vee\}$ 与 $\{\oplus,\neg\}$ 不是全功能集。}{\textbf{Completeness}: $\{\neg,\wedge,\vee\}$ is functionally complete; by De Morgan, $\{\neg,\wedge\}$, $\{\neg,\vee\}$ and $\{\neg,\to\}$ are also complete; the singletons $\{\uparrow\}$ and $\{\downarrow\}$ are each functionally complete. Conversely, $\{\wedge,\vee\}$ and $\{\oplus,\neg\}$ are not functionally complete.}
\PRF{由范式存在定理，任何公式等价于一个 DNF，而 DNF 只用 $\neg,\wedge,\vee$，故 $\{\neg,\wedge,\vee\}$ 完备。$\uparrow$ 的情形：$\neg P\Leftrightarrow P\uparrow P$，进而 $P\wedge Q\Leftrightarrow(P\uparrow Q)\uparrow(P\uparrow Q)$、$P\vee Q\Leftrightarrow(P\uparrow P)\uparrow(Q\uparrow Q)$，再借 $\{\neg,\wedge,\vee\}$ 完备即得。不完整性的判定依据是“封闭性质”：$\wedge,\vee$ 是单调函数，复合仍是单调的，而 $\neg$ 非单调，故 $\{\wedge,\vee\}$ 表示不了 $\neg$；$\oplus$ 与 $\neg$ 都是线性（仿射）函数，复合仍是线性的，而 $P\wedge Q$ 非线性，故 $\{\oplus,\neg\}$ 表示不了 $\wedge$。}{By the normal-form theorem, every formula is equivalent to a DNF, which uses only $\neg,\wedge,\vee$; hence $\{\neg,\wedge,\vee\}$ is complete. For $\uparrow$: $\neg P\Leftrightarrow P\uparrow P$, then $P\wedge Q\Leftrightarrow(P\uparrow Q)\uparrow(P\uparrow Q)$ and $P\vee Q\Leftrightarrow(P\uparrow P)\uparrow(Q\uparrow Q)$, and completeness follows from that of $\{\neg,\wedge,\vee\}$. Incompleteness is shown by a closedness property: $\wedge,\vee$ are monotone and so are their composites, while $\neg$ is not monotone, so $\{\wedge,\vee\}$ cannot express $\neg$; $\oplus$ and $\neg$ are both linear (affine) and so are their composites, while $P\wedge Q$ is not linear, so $\{\oplus,\neg\}$ cannot express $\wedge$.}
\REM{更一般地，Post 定理刻画了全部全功能集：联结词集是完备的当且仅当它不整体落在五个“极大封闭类”（保真、保假、单调、仿射、自对偶）中的任何一个。}{More generally, Post's theorem characterizes all functionally complete sets: a set of connectives is complete iff it is not contained in any of the five maximal clones (truth-preserving, falsity-preserving, monotone, affine/linear, self-dual).}

\subsection{推理规则与证明方法 | Inference rules and proof methods}
\DEF{推理形式“由前提 $A_1,\dots,A_k$ 推出结论 $B$”记 $A_1\wedge\cdots\wedge A_k\Rightarrow B$，并称其为\textbf{有效推理}当且仅当 $A_1\wedge\cdots\wedge A_k\to B$ 是重言式（此时也叫\textbf{可靠}，sound）。}{The inference form ``from premises $A_1,\dots,A_k$ infer $B$'' is written $A_1\wedge\cdots\wedge A_k\Rightarrow B$ and is \textbf{valid} (sound) exactly when $A_1\wedge\cdots\wedge A_k\to B$ is a tautology.}
\THM{\textbf{演绎定理（CP 规则）}：$A_1,\dots,A_k\Rightarrow A\to B$ 当且仅当 $A_1,\dots,A_k,A\Rightarrow B$。常用推理规则：假言推理、拒取式、假言三段论 $((P\to Q)\wedge(Q\to R))\Rightarrow(P\to R)$、析取三段论 $((P\vee Q)\wedge\neg P)\Rightarrow Q$、附加、化简、合取。\textbf{归结原理}：把前提与结论的否定化为子句集（子句是文字的析取），反复应用消解规则 $\{A\vee C,\ \neg A\vee D\}\rightsquigarrow\{C\vee D\}$；归结是\textbf{反驳完备}的——子句集不可满足当且仅当可由它消解出空子句 $\square$。}{\textbf{Deduction theorem (CP rule)}: $A_1,\dots,A_k\Rightarrow A\to B$ iff $A_1,\dots,A_k,A\Rightarrow B$. Usual rules: modus ponens, modus tollens, hypothetical syllogism $((P\to Q)\wedge(Q\to R))\Rightarrow(P\to R)$, disjunctive syllogism $((P\vee Q)\wedge\neg P)\Rightarrow Q$, addition, simplification and conjunction. \textbf{Resolution}: convert the premises and the negation of the conclusion to a clause set (a clause is a disjunction of literals) and repeatedly apply the resolution rule $\{A\vee C,\ \neg A\vee D\}\rightsquigarrow\{C\vee D\}$. Resolution is \textbf{refutation complete}: a clause set is unsatisfiable iff the empty clause $\square$ can be derived from it.}
\PRF{演绎定理：把 $A_1\wedge\cdots\wedge A_k\to(A\to B)$ 用等价式展开，与 $(A_1\wedge\cdots\wedge A_k\wedge A)\to B$ 相等。归结完备性的核心是“消解保持可满足性”：若 $\{A\vee C,\neg A\vee D\}$ 在某个赋值下都为真，则无论 $A$ 真假，$C\vee D$ 也必为真；反复消解直到不能再消解，若未得 $\square$，可给剩余文字赋值使其全部为真，从而原集可满足。三种证明方法对应：\textbf{直接证明}（前提出发推结论）、\textbf{附加前提法}（CP，把 $A\to B$ 转为加 $A$ 证 $B$）、\textbf{反证法}（加 $\neg B$ 推出矛盾 $C\wedge\neg C$）。}{Deduction theorem: expand $A_1\wedge\cdots\wedge A_k\to(A\to B)$ by equivalences; it equals $(A_1\wedge\cdots\wedge A_k\wedge A)\to B$. The core of resolution completeness is that resolution preserves satisfiability: if $\{A\vee C,\neg A\vee D\}$ both hold under an assignment, then $C\vee D$ holds regardless of $A$; repeat until no resolution is possible, and if $\square$ is not obtained, assign the remaining literals to make them true, so the original set is satisfiable. The three proof methods correspond to \textbf{direct proof} (derive the conclusion from premises), \textbf{conditional proof} (CP: turn $A\to B$ into adding $A$ and proving $B$) and \textbf{indirect proof} (add $\neg B$ and derive a contradiction $C\wedge\neg C$).}

\subsection{谓词、量词与辖域 | Predicates, quantifiers and scope}
\DEF{\textbf{谓词} $P(x_1,\dots,x_n)$ 是带 $n$ 个空位的断言，代入个体后成为命题；个体变元的取值范围叫\textbf{个体域}（论述域）。\textbf{全称量词} $\forall x$ 表“对个体域中所有 $x$”，\textbf{存在量词} $\exists x$ 表“存在 $x$”，$\exists!x$ 表“存在唯一 $x$”。量词后紧接的子公式叫它的\textbf{辖域}；辖域内被该量词约束的变元出现叫\textbf{约束出现}，否则叫\textbf{自由出现}。不含自由变元的公式叫\textbf{闭式}（语句）。自由变元可代入，约束变元可换名（$\alpha$ 改名），但不得改变约束关系或造成变元被捕获。}{A \textbf{predicate} $P(x_1,\dots,x_n)$ is an assertion with $n$ blanks that becomes a proposition once individuals are substituted; the range of the individual variables is the \textbf{domain} (universe of discourse). The \textbf{universal quantifier} $\forall x$ means ``for all $x$ in the domain'', the \textbf{existential quantifier} $\exists x$ means ``there exists $x$'', and $\exists!x$ ``there exists a unique $x$''. The subformula immediately following a quantifier is its \textbf{scope}; occurrences of a variable bound by a quantifier are \textbf{bound}, the others \textbf{free}. A formula without free variables is \textbf{closed} (a sentence). Free variables may be substituted and bound variables may be renamed ($\alpha$-renaming), provided the binding structure is preserved and no variable is captured.}
\REM{翻译口诀：\textbf{全称句用蕴含，存在句用合取}。若个体域已是“人”，则“所有人都会死”译为 $\forall x\,\mathrm{Mortal}(x)$；若为全总域，则写 $\forall x(\mathrm{Man}(x)\to\mathrm{Mortal}(x))$。存在句“有人会飞”是 $\exists x(\mathrm{Man}(x)\wedge\mathrm{Fly}(x))$，不能写 $\exists x(\mathrm{Man}(x)\to\mathrm{Fly}(x))$，否则只要有一个非人个体它就为真。}{Translation rule of thumb: \textbf{universal statements take implication, existential statements take conjunction}. If the domain is already ``people'', ``all people are mortal'' is $\forall x\,\mathrm{Mortal}(x)$; over a universal domain it is $\forall x(\mathrm{Man}(x)\to\mathrm{Mortal}(x))$. The existential ``someone can fly'' is $\exists x(\mathrm{Man}(x)\wedge\mathrm{Fly}(x))$, not $\exists x(\mathrm{Man}(x)\to\mathrm{Fly}(x))$, which would be true as soon as one non-human individual exists.}

\subsection{谓词演算的永真公式与推理规则 | Valid formulas and inference rules of predicate calculus}
\DEF{谓词公式在\textbf{解释}（一个个体域、对谓词与自由变元的指派）下取真值；在所有解释下都真才是\textbf{永真公式}，在某个解释下为真则\textbf{可满足}。}{A predicate formula takes a truth value under an \textbf{interpretation} (a domain together with assignments to predicates and free variables); it is \textbf{valid} when true under every interpretation and \textbf{satisfiable} when true under some interpretation.}
\THM{\textbf{量词的基本等价与永真式}（其中 $x$ 不在 $P$ 中自由出现）：量词否定律 $\neg\forall xP(x)\Leftrightarrow\exists x\neg P(x)$、$\neg\exists xP(x)\Leftrightarrow\forall x\neg P(x)$；分配律 $\forall x(P\wedge Q)\Leftrightarrow\forall xP\wedge\forall xQ$、$\exists x(P\vee Q)\Leftrightarrow\exists xP\vee\exists xQ$；空量词 $\forall xP\Leftrightarrow P$、$\exists xP\Leftrightarrow P$；辖域收缩 $\forall x(P\to Q(x))\Leftrightarrow(P\to\forall xQ(x))$、$\exists x(P\wedge Q(x))\Leftrightarrow(P\wedge\exists xQ(x))$；以及单向蕴含 $\exists y\forall xP(x,y)\Rightarrow\forall x\exists yP(x,y)$，其逆不成立。}{Let $x$ not occur free in $P$. Quantifier negation: $\neg\forall xP(x)\Leftrightarrow\exists x\neg P(x)$ and $\neg\exists xP(x)\Leftrightarrow\forall x\neg P(x)$. Distribution: $\forall x(P\wedge Q)\Leftrightarrow\forall xP\wedge\forall xQ$ and $\exists x(P\vee Q)\Leftrightarrow\exists xP\vee\exists xQ$. Null quantification: $\forall xP\Leftrightarrow P$ and $\exists xP\Leftrightarrow P$. Scope reduction: $\forall x(P\to Q(x))\Leftrightarrow(P\to\forall xQ(x))$ and $\exists x(P\wedge Q(x))\Leftrightarrow(P\wedge\exists xQ(x))$. One-way implication: $\exists y\forall xP(x,y)\Rightarrow\forall x\exists yP(x,y)$, whose converse fails.}
\THM{\textbf{推理规则与限制}：全称实例化 UI $\forall xP(x)\vdash P(c)$；全称推广 UG，若对任意个体 $c$ 证得 $P(c)$ 则 $\forall xP(x)$；存在实例化 EI $\exists xP(x)\vdash P(c)$，其中 $c$ 必须是此前未出现的新常元；存在推广 EG $P(c)\vdash\exists xP(x)$。关键限制：由 EI 引入的常元不能再做 UG，否则会把“存在”偷换成“任意”，从而从 $\exists xP(x)$ 错误推出 $\forall xP(x)$。}{Four inference rules with restrictions: universal instantiation UI, $\forall xP(x)\vdash P(c)$; universal generalization UG, from $P(c)$ for an arbitrary $c$ infer $\forall xP(x)$; existential instantiation EI, $\exists xP(x)\vdash P(c)$ with $c$ a new constant not used before; existential generalization EG, $P(c)\vdash\exists xP(x)$. The crucial restriction: a constant introduced by EI may not later be generalized by UG, otherwise ``there exists'' is smuggled into ``for all'', incorrectly inferring $\forall xP(x)$ from $\exists xP(x)$.}

\EX{判定 $((P\to Q)\wedge P)\to Q$ 是否永真：列 $P,Q$ 的 $2^2=4$ 种赋值，$P\to Q$ 只在 $(1,0)$ 行为假，故 $(P\to Q)\wedge P$ 只在第一行为真，该行 $Q$ 为真，蕴含式为真；其余三行前件为假，蕴含式仍为真。四行全为 $1$，故它是重言式。真值表给出判定永真性的机械算法，代价是 $2^n$ 行。}{To decide whether $((P\to Q)\wedge P)\to Q$ is a tautology, list the $2^2=4$ assignments. The column $P\to Q$ is false only in row $(1,0)$, so $(P\to Q)\wedge P$ is true only in the first row, where $Q$ is true and the implication is true; in the other three rows the antecedent is false and the implication is again true. All four rows are $1$, so the formula is a tautology. Truth tables give a mechanical decision procedure for validity at the cost of $2^n$ rows.}
\begin{center}\small
\begin{tabular}{@{}cccccc@{}}
\toprule
$P$ & $Q$ & $P\to Q$ & $(P\to Q)\wedge P$ & $((P\to Q)\wedge P)\to Q$ \\
\midrule
1 & 1 & 1 & 1 & 1 \\
1 & 0 & 0 & 0 & 1 \\
0 & 1 & 1 & 0 & 1 \\
0 & 0 & 1 & 0 & 1 \\
\bottomrule
\end{tabular}
\end{center}
\EX{$\forall x\exists y$ 与 $\exists y\forall x$ 不同：个体域取整数、$P(x,y)$ 为 $x<y$ 时，$\forall x\exists y(x<y)$ 真（每个整数都有更大的整数），而 $\exists y\forall x(x<y)$ 假（不存在大于一切整数的整数）。这正说明 $\exists y\forall xP\Rightarrow\forall x\exists yP$ 的逆不成立。}{$\forall x\exists y$ differs from $\exists y\forall x$: on the integers with $P(x,y)$ meaning $x<y$, the formula $\forall x\exists y(x<y)$ is true (every integer has a larger one) while $\exists y\forall x(x<y)$ is false (no integer exceeds every integer). This shows the converse of $\exists y\forall xP\Rightarrow\forall x\exists yP$ fails.}

\ASK{为什么 $P\to Q$ 在前件为假时规定为真？}{Why is $P\to Q$ declared true when $P$ is false?}{实质蕴含只排除“前件真、后件假”，这样 $P\to Q\Leftrightarrow\neg P\vee Q$，假言推理、拒取式才在所有赋值下可靠，逆反命题才与原命题等值；若改用因果蕴含，推理规则就无法用真值表验证。}{Material implication excludes only ``antecedent true, consequent false''; this gives $P\to Q\Leftrightarrow\neg P\vee Q$, keeps modus ponens and modus tollens sound under all assignments, and makes the contrapositive equivalent to the original. A causal implication could not be checked by truth tables.}
\ASK{主范式为什么唯一？}{Why is the principal normal form unique?}{公式的真值表是唯一确定的；把为真的行对应的极小项全部析取、为假的行对应的极大项全部合取，两种主范式由真值表唯一决定，因此在不计项与文字的次序及重复的意义下唯一。DNF/CNF 一般不唯一，主范式是其中的规范形式。}{A formula's truth table is uniquely determined; the disjunction of the minterms of true rows and the conjunction of the maxterms of false rows are fixed by the truth table, so the principal forms are unique up to order and repetition. DNF/CNF are generally not unique; the principal forms are their canonical representatives.}

\section[二元关系（第 3 章）]{二元关系（教材第 3 章）| Binary relations (Chapter 3)}
\subsection{关系的定义、表示与性质 | Definition, representations and properties}
\DEF{从集合 $A$ 到集合 $B$ 的\textbf{二元关系}是笛卡儿积的子集 $R\subseteq A\times B$；$(a,b)\in R$ 也写作 $aRb$。$A=B$ 时称 $A$ 上的关系。有限关系有三种等价表示：序偶集合、\textbf{关系矩阵} $M_R=(r_{ij})$（$a_iRb_j$ 时 $r_{ij}=1$，否则 $0$）、\textbf{关系图}（顶点为元素、有向边表示序偶）。}{A \textbf{binary relation} from $A$ to $B$ is a subset $R\subseteq A\times B$; $(a,b)\in R$ is also written $aRb$. When $A=B$ one speaks of a relation on $A$. A finite relation has three equivalent representations: a set of ordered pairs, a \textbf{relation matrix} $M_R=(r_{ij})$ with $r_{ij}=1$ iff $a_iRb_j$, and a \textbf{digraph} with vertices the elements and directed edges the pairs.}
\DEF{$A$ 上的关系 $R$ 的五个基本性质：\textbf{自反} $\forall a\,aRa$；\textbf{反自反} $\forall a\,\neg aRa$；\textbf{对称} $aRb\Rightarrow bRa$；\textbf{反对称} $aRb\wedge bRa\Rightarrow a=b$；\textbf{传递} $aRb\wedge bRc\Rightarrow aRc$。注意：反自反不是自反的否定，反对称也不是对称的否定。}{The five basic properties of a relation $R$ on $A$: \textbf{reflexive} $\forall a\,aRa$; \textbf{irreflexive} $\forall a\,\neg aRa$; \textbf{symmetric} $aRb\Rightarrow bRa$; \textbf{antisymmetric} $aRb\wedge bRa\Rightarrow a=b$; \textbf{transitive} $aRb\wedge bRc\Rightarrow aRc$. Note: irreflexive is not the negation of reflexive, nor antisymmetric the negation of symmetric.}
\THM{\textbf{矩阵刻画}：$R$ 自反 $\Leftrightarrow$ $M_R$ 主对角线全为 $1$；反自反 $\Leftrightarrow$ 主对角线全为 $0$；对称 $\Leftrightarrow$ $M_R$ 是对称矩阵；反对称 $\Leftrightarrow$ 对一切 $i\ne j$，$r_{ij}$ 与 $r_{ji}$ 不同时为 $1$；传递 $\Leftrightarrow$ $M_R^2\le M_R$（布尔意义，即 $(M_R\odot M_R)_{ij}=1$ 时必有 $r_{ij}=1$）。}{In matrices: $R$ is reflexive iff the main diagonal of $M_R$ is all $1$; irreflexive iff it is all $0$; symmetric iff $M_R$ is symmetric; antisymmetric iff for all $i\ne j$, $r_{ij}$ and $r_{ji}$ are not both $1$; transitive iff $M_R^2\le M_R$ in the Boolean sense, i.e.\ $(M_R\odot M_R)_{ij}=1$ implies $r_{ij}=1$.}

\subsection{关系的合成 | Composition of relations}
\DEF{设 $R\subseteq A\times B$、$S\subseteq B\times C$，合成关系 $R\circ S\subseteq A\times C$ 定义为
\[
R\circ S=\{(a,c)\mid \exists b\in B,\ (a,b)\in R\ \wedge\ (b,c)\in S\}.
\]
本教材采用\textbf{从左到右}的次序（先 $R$ 后 $S$），与函数复合 $g\circ f$（先 $f$ 后 $g$）方向相反，读题时须确认记号。$R$ 的幂递归定义为 $R^1=R$，$R^{n+1}=R^n\circ R$。}{Let $R\subseteq A\times B$ and $S\subseteq B\times C$; the composite $R\circ S\subseteq A\times C$ is defined by the displayed formula. This textbook uses the \textbf{left-to-right} order (first $R$, then $S$), opposite to function composition $g\circ f$ (first $f$, then $g$), so the convention must be checked. Powers are defined by $R^1=R$ and $R^{n+1}=R^n\circ R$.}
\THM{\textbf{合成性质}：关系合成满足结合律 $(R\circ S)\circ T=R\circ(S\circ T)$，一般不满足交换律；矩阵上 $M_{R\circ S}=M_R\odot M_S$，其中 $\odot$ 是布尔乘法（把普通加、乘换成 $\vee,\wedge$）。于是 $M_{R^n}=(M_R)^n$（布尔幂）。}{\textbf{Composition}: composition of relations is associative, $(R\circ S)\circ T=R\circ(S\circ T)$, but in general not commutative; in matrices $M_{R\circ S}=M_R\odot M_S$, where $\odot$ is Boolean multiplication (replace $+,\cdot$ by $\vee,\wedge$). Hence $M_{R^n}=(M_R)^n$ (Boolean power).}
\PRF{结合律由定义逐点验证：$(a,c)\in(R\circ S)\circ T$ 与 $(a,c)\in R\circ(S\circ T)$ 都等价于“存在 $b,d$ 使 $(a,b)\in R$、$(b,d)\in S$、$(d,c)\in T$”。矩阵式是矩阵乘法的直接翻译：$(M_R\odot M_S)_{ik}=\bigvee_j(r_{ij}\wedge s_{jk})$。}{Associativity is verified pointwise: $(a,c)\in(R\circ S)\circ T$ and $(a,c)\in R\circ(S\circ T)$ are both equivalent to ``there exist $b,d$ with $(a,b)\in R$, $(b,d)\in S$, $(d,c)\in T$''. The matrix formula translates matrix multiplication directly: $(M_R\odot M_S)_{ik}=\bigvee_j(r_{ij}\wedge s_{jk})$.}

\subsection{关系的闭包 | Closures of relations}
\DEF{设 $R$ 是 $A$ 上的关系，$P$ 为某一性质（自反、对称或传递）。$R$ 的\textbf{$P$ 闭包}是包含 $R$ 且具有性质 $P$ 的\textbf{最小}关系；分别记为自反闭包 $r(R)$、对称闭包 $s(R)$、传递闭包 $t(R)$。等价地，$P$ 闭包等于所有包含 $R$ 且具有性质 $P$ 的关系的交。}{Let $R$ be a relation on $A$ and $P$ a property (reflexive, symmetric or transitive). The \textbf{$P$-closure} of $R$ is the \textbf{smallest} relation containing $R$ with property $P$; they are written reflexive closure $r(R)$, symmetric closure $s(R)$ and transitive closure $t(R)$. Equivalently, the $P$-closure is the intersection of all relations containing $R$ with property $P$.}
\THM{\textbf{闭包构造}：设 $|A|=n$，$I_A=\{(a,a)\mid a\in A\}$，$R^{-1}=\{(b,a)\mid aRb\}$，则
\[
r(R)=R\cup I_A,\qquad s(R)=R\cup R^{-1},\qquad t(R)=R\cup R^2\cup\cdots\cup R^n .
\]
}{Let $|A|=n$, $I_A=\{(a,a)\mid a\in A\}$ and $R^{-1}=\{(b,a)\mid aRb\}$. Then the three displayed formulas hold.}
\PRF{$r(R),s(R)$ 显然包含 $R$ 且具有相应性质，并含于任何更大的同类关系。$t(R)$ 的关键是“只需并到 $R^n$”：$R^k$ 对应长度 $k$ 的有向路径，若 $k>n$，路径经过 $n+1$ 个顶点，由抽屉原理必重复某顶点，删去中间回路得到更短路径，故该序偶已出现在某个 $R^j$（$j\le n$）中。}{$r(R)$ and $s(R)$ clearly contain $R$, have the required property, and are contained in any larger relation with it. The key point for $t(R)$ is ``powers up to $R^n$ suffice'': $R^k$ corresponds to directed walks of length $k$, and if $k>n$ the walk visits $n+1$ vertices, so by the pigeonhole principle some vertex repeats; deleting the cycle between the two occurrences leaves a shorter walk, so the pair already lies in some $R^j$ with $j\le n$.}
\THM{\textbf{Warshall 算法}：令 $W^{(0)}=M_R$，对 $k=1,\dots,n$ 计算
\[
w_{ij}^{(k)}=w_{ij}^{(k-1)}\vee\bigl(w_{ik}^{(k-1)}\wedge w_{kj}^{(k-1)}\bigr),
\]
则 $W^{(n)}=M_{t(R)}$；其含义是“允许 $1,\dots,k$ 号顶点作中转”，时间复杂度 $O(n^3)$。}{Start from $W^{(0)}=M_R$ and for $k=1,\dots,n$ compute the displayed recurrence; then $W^{(n)}=M_{t(R)}$. Its meaning is ``vertices $1,\dots,k$ are allowed as intermediate stops''; it runs in $O(n^3)$.}
\PRF{归纳地，$W^{(k)}$ 记录“中间顶点取自 $\{1,\dots,k\}$”的可达性：新可达路径要么不用 $k$（$w_{ij}^{(k-1)}$），要么用 $k$ 一次（$w_{ik}^{(k-1)}\wedge w_{kj}^{(k-1)}$）；$k=n$ 时中间顶点可取全部，即传递闭包。}{By induction, $W^{(k)}$ records reachability with intermediate vertices from $\{1,\dots,k\}$: a new reachable path either avoids $k$ ($w_{ij}^{(k-1)}$) or uses $k$ once ($w_{ik}^{(k-1)}\wedge w_{kj}^{(k-1)}$); when $k=n$ all vertices are allowed, which is the transitive closure.}

\subsection{次序关系 | Order relations}
\DEF{$A$ 上的关系若自反、反对称、传递，称为\textbf{偏序关系}，记作 $\preccurlyeq$；$(A,\preccurlyeq)$ 叫\textbf{偏序集}。若偏序还满足任意两元素可比较（$a\preccurlyeq b$ 或 $b\preccurlyeq a$），则称\textbf{全序}（线序）。典型：实数上的 $\le$ 是全序；$\subseteq$ 是偏序但一般非全序；正整数上的整除是偏序。}{A relation on $A$ that is reflexive, antisymmetric and transitive is a \textbf{partial order}, written $\preccurlyeq$, and $(A,\preccurlyeq)$ is a \textbf{partially ordered set}. If in addition every two elements are comparable ($a\preccurlyeq b$ or $b\preccurlyeq a$), it is a \textbf{total (linear) order}. Typical: $\le$ on the reals is total; $\subseteq$ is partial but generally not total; divisibility on the positive integers is partial.}
\DEF{\textbf{哈斯图}只画\textbf{覆盖}关系：$b$ 覆盖 $a$ 指 $a\prec b$（$a\preccurlyeq b$ 且 $a\ne b$）且不存在 $c$ 使 $a\prec c\prec b$；边向上画、不画箭头。设 $B\subseteq A$。$A$ 中元素 $m$ 是 $B$ 的\textbf{上界}若 $\forall b\in B,\ b\preccurlyeq m$；$B$ 的\textbf{上确界}（最小上界，lub，记 $\vee B$）是最小的上界；下界与\textbf{下确界}（最大下界，glb，记 $\wedge B$）对偶。$A$ 的\textbf{极大元}是上方无元素的元素，\textbf{最大元}是大于等于所有元素的元素；极小元、最小元对偶。}{A \textbf{Hasse diagram} draws only the \textbf{cover} relation: $b$ covers $a$ when $a\prec b$ ($a\preccurlyeq b$ and $a\ne b$) and no $c$ satisfies $a\prec c\prec b$; edges point upward without arrows. Let $B\subseteq A$. An element $m$ of $A$ is an \textbf{upper bound} of $B$ if $b\preccurlyeq m$ for all $b\in B$; the \textbf{least upper bound} (supremum, lub, written $\vee B$) is the least upper bound; lower bounds and the \textbf{greatest lower bound} (infimum, glb, written $\wedge B$) are dual. A \textbf{maximal element} of $A$ has nothing above it and a \textbf{greatest element} is above or equal to every element; minimal and least are dual.}
\THM{\textbf{确界的唯一性}：最大元、最小元、上确界、下确界若存在必唯一；但极大元、极小元可以有多个。最大元必为唯一极大元。在有限偏序集中，每个元素都在某个极大元之下，因此唯一极大元也是最大元。无限偏序集可取一条无上端的自然数链，再加入与链上所有元素不可比的单点。这个单点是唯一极大元，但并非最大元。}{If they exist, the greatest element, least element, least upper bound and greatest lower bound are unique; maximal and minimal elements may be many. A greatest element is the unique maximal element. In a finite poset, every element lies below some maximal element, so uniqueness implies greatestness. The disjoint union of the increasing natural-number chain and an incomparable singleton shows why finiteness matters.}
\PRF{唯一性来自反对称性：若 $m,m'$ 都是上确界，则 $m\preccurlyeq m'$ 且 $m'\preccurlyeq m$，故 $m=m'$。}{Uniqueness follows from antisymmetry: if $m,m'$ are both least upper bounds, then $m\preccurlyeq m'$ and $m'\preccurlyeq m$, hence $m=m'$.}

\subsection{等价关系、等价类与划分 | Equivalence relations, classes and partitions}
\DEF{$A$ 上的关系若自反、对称、传递，称为\textbf{等价关系}，记 $\sim$。$a$ 的\textbf{等价类} $[a]=\{x\in A\mid x\sim a\}$；全体等价类组成的集合叫\textbf{商集} $A/{\sim}$。$A$ 的\textbf{划分}是 $A$ 的一族非空子集（块），两两不交且并为 $A$。}{A relation on $A$ that is reflexive, symmetric and transitive is an \textbf{equivalence relation}, written $\sim$. The \textbf{equivalence class} of $a$ is $[a]=\{x\in A\mid x\sim a\}$; the set of all classes is the \textbf{quotient set} $A/{\sim}$. A \textbf{partition} of $A$ is a family of nonempty subsets (blocks), pairwise disjoint with union $A$.}
\THM{\textbf{等价类的基本性质}：$a\in[a]$；$[a]=[b]\Leftrightarrow a\sim b$；$[a]\ne[b]\Rightarrow[a]\cap[b]=\varnothing$。故等价类两两不交且并为 $A$，构成 $A$ 的一个划分。}{Let $\sim$ be an equivalence relation on $A$. Then $a\in[a]$; $[a]=[b]\Leftrightarrow a\sim b$; and if $[a]\ne[b]$ then $[a]\cap[b]=\varnothing$. Hence the classes are pairwise disjoint with union $A$ and form a partition of $A$.}
\THM{\textbf{等价关系与划分的对应}：$A$ 上的等价关系与 $A$ 的划分之间存在一一对应：等价关系对应其商集（块即等价类）；划分对应关系 $a\sim b\Leftrightarrow a,b$ 属于同一块。}{\textbf{Correspondence}: there is a bijection between equivalence relations on $A$ and partitions of $A$: an equivalence relation maps to its quotient set (blocks are the classes), and a partition maps to the relation $a\sim b\Leftrightarrow a,b$ lie in the same block.}
\PRF{第一定理由对称与传递直接验证：若 $x\in[a]\cap[b]$，则 $x\sim a$、$x\sim b$，由对称得 $a\sim x$，由传递得 $a\sim b$，再证两集合相等。第二定理需证两个方向互为逆：由等价关系到划分再到“同块”关系，恢复出原关系（因 $a\sim b\Leftrightarrow[a]=[b]$）；由划分到关系再到商集，恢复出原划分。}{For the first theorem, symmetry and transitivity give: if $x\in[a]\cap[b]$, then $x\sim a$ and $x\sim b$, so $a\sim x$ and transitivity yields $a\sim b$, and then the sets are equal. For the second, one verifies the two maps are mutually inverse: an equivalence relation, mapped to its partition and back to ``same block'', recovers the original relation (since $a\sim b\Leftrightarrow[a]=[b]$); a partition, mapped to a relation and back to its quotient set, recovers the original partition.}

\EX{在整数集上定义 $a\equiv b\pmod n$ 当且仅当 $n\mid(a-b)$。这是等价关系，等价类为 $[k]=\{k+tn\mid t\in\mathbb Z\}$（$k=0,\dots,n-1$），商集记 $\mathbb Z_n$；$[a]=[b]\Leftrightarrow a\equiv b\pmod n$。等价类是模 $n$ 的剩余类，其全体构成 $\mathbb Z$ 的划分。}{On the integers define $a\equiv b\pmod n$ iff $n\mid(a-b)$. This is an equivalence relation with classes $[k]=\{k+tn\mid t\in\mathbb Z\}$ for $k=0,\dots,n-1$; the quotient set is written $\mathbb Z_n$, and $[a]=[b]\Leftrightarrow a\equiv b\pmod n$. The classes are the residue classes modulo $n$, and together they partition $\mathbb Z$.}

\ASK{极大元与最大元、上界与上确界有何区别？}{What distinguishes maximal from greatest elements, and upper bounds from least upper bounds?}{最大元必须与所有元素可比且大于等于它们，存在则唯一；极大元只要求上方无元素，可以多个。上界是“上方元素”（可不在子集内），上确界是上界中最小的；按教材第 3.4 节，上确界也叫最小上界，记作 $\operatorname{lub}$。即使偏序集有限，上界也可能有多个彼此不可比的极小元素。比如令 $a,b$ 都严格小于 $c,d$，除此之外元素互不可比。则 $\{a,b\}$ 的上界有 $c,d$，却无最小上界。}{A greatest element must be comparable to and above every element, and is unique if it exists; a maximal element only needs nothing above it, and there may be several. An upper bound is an element above the subset (not necessarily in it), and the least upper bound is the least of all upper bounds; Even a finite poset may have upper bounds without a least one. If $a,b$ are both below incomparable maximal elements $c,d$ and there are no other strict comparisons, then $c,d$ are upper bounds of $\{a,b\}$ but neither is least.}
\ASK{为什么传递闭包只需并到 $R^n$？}{Why does the transitive closure stop at $R^n$?}{$R^k$ 对应长度 $k$ 的有向路径；$k>n$ 时路径经过 $n+1$ 个顶点，抽屉原理保证有重复顶点，删去中间回路得更短路径，故该序偶已属于某个 $R^j$（$j\le n$）。因此 $t(R)=R\cup\cdots\cup R^n$，Warshall 算法也只需 $n$ 轮中转。}{$R^k$ corresponds to directed walks of length $k$; if $k>n$ a walk visits $n+1$ vertices, so by the pigeonhole principle some vertex repeats, and deleting the intermediate cycle leaves a shorter walk, so the pair already belongs to some $R^j$ with $j\le n$. Hence $t(R)=R\cup\cdots\cup R^n$, and Warshall's algorithm needs only $n$ rounds of intermediate vertices.}

\section[代数（第 6 章）]{代数（教材第 6 章）| Algebra (Chapter 6)}
\subsection{代数结构、子代数与同态 | Algebraic structures, subalgebras and homomorphisms}
\DEF{\textbf{代数结构}（代数系统）由非空集合 $A$（\textbf{载体}）、若干 $A$ 上的\textbf{运算}（$A^k\to A$ 的映射，$k$ 为元数）及若干常元组成；运算必须对 $A$ \textbf{封闭}。\textbf{子代数}是载体的子集，对一切运算封闭且含全部常元。两个同类型代数 $(A,*)$、$(B,\oplus)$ 之间的映射 $h:A\to B$ 若满足 $h(a*b)=h(a)\oplus h(b)$（对每个运算，并保持常元），称为\textbf{同态}；单射、满射、双射的同态分别叫\textbf{单同态}、\textbf{满同态}、\textbf{同构}（记 $\cong$）；$A$ 到自身的同态叫自同态，双射自同态叫自同构。}{An \textbf{algebraic structure} consists of a nonempty set $A$ (the \textbf{carrier}), several \textbf{operations} on $A$ (maps $A^k\to A$, $k$ the arity) and distinguished constants; operations must be \textbf{closed} on $A$. A \textbf{subalgebra} is a subset of the carrier closed under all operations and containing all constants. A map $h:A\to B$ between algebras $(A,*)$, $(B,\oplus)$ of the same type with $h(a*b)=h(a)\oplus h(b)$ for every operation (and preserving constants) is a \textbf{homomorphism}; injective, surjective and bijective ones are \textbf{monomorphism}, \textbf{epimorphism} and \textbf{isomorphism} (written $\cong$); an endomorphism is a homomorphism $A\to A$, and a bijective endomorphism is an automorphism.}
\DEF{设 $h:A\to B$ 是同态，$e_B$ 为 $B$ 的单位元。$h$ 的\textbf{核} $\ker h=\{a\in A\mid h(a)=e_B\}$，\textbf{像} $\operatorname{Im}h=\{h(a)\mid a\in A\}$。}{For a homomorphism $h:A\to B$ with $e_B$ the identity of $B$, the \textbf{kernel} is $\ker h=\{a\in A\mid h(a)=e_B\}$ and the \textbf{image} is $\operatorname{Im}h=\{h(a)\mid a\in A\}$.}
\THM{\textbf{同态的基本性质}（群情形）：$h(e_A)=e_B$，$h(a^{-1})=h(a)^{-1}$；$\operatorname{Im}h$ 是 $B$ 的子群，$\ker h$ 是 $A$ 的正规子群；$h$ 是单同态当且仅当 $\ker h=\{e_A\}$。}{For groups: $h(e_A)=e_B$ and $h(a^{-1})=h(a)^{-1}$; $\operatorname{Im}h$ is a subgroup of $B$ and $\ker h$ is a normal subgroup of $A$; $h$ is injective iff $\ker h=\{e_A\}$.}
\PRF{由 $h(e_A)=h(e_Ae_A)=h(e_A)h(e_A)$ 左乘逆元得 $h(e_A)=e_B$；由 $h(a)h(a^{-1})=h(aa^{-1})=h(e_A)=e_B$ 得 $h(a^{-1})=h(a)^{-1}$。单射性：若 $h(a)=h(b)$，则 $h(ab^{-1})=e_B$，即 $ab^{-1}\in\ker h$；核平凡时 $ab^{-1}=e_A$，故 $a=b$。}{Since $h(e_A)=h(e_Ae_A)=h(e_A)h(e_A)$, left-multiplying by the inverse gives $h(e_A)=e_B$; from $h(a)h(a^{-1})=h(aa^{-1})=h(e_A)=e_B$ we get $h(a^{-1})=h(a)^{-1}$. For injectivity: if $h(a)=h(b)$, then $h(ab^{-1})=e_B$, i.e.\ $ab^{-1}\in\ker h$; with a trivial kernel $ab^{-1}=e_A$, so $a=b$.}

\subsection{同余关系、商代数与积代数 | Congruences, quotient algebras and products}
\DEF{代数 $(A,*)$ 上的等价关系 $\sim$ 若与运算\textbf{相容}——$a\sim b$ 且 $c\sim d$ 蕴含 $a*c\sim b*d$（对每个运算）——则称为\textbf{同余关系}。在商集 $A/{\sim}$ 上定义 $[a]*[b]=[a*b]$，所得代数叫\textbf{商代数}。两个同类型代数的\textbf{积代数}载体为 $A\times B$，运算按分量进行。}{An equivalence relation $\sim$ on $(A,*)$ is a \textbf{congruence} if it is \textbf{compatible} with the operations: $a\sim b$ and $c\sim d$ imply $a*c\sim b*d$ for every operation. Defining $[a]*[b]=[a*b]$ on $A/{\sim}$ gives the \textbf{quotient algebra}. The \textbf{product algebra} of two algebras of the same type has carrier $A\times B$ and componentwise operations.}
\THM{\textbf{商代数的良定义性}：同余关系 $\sim$ 使 $[a]*[b]=[a*b]$ 与代表元选取无关，即商代数的运算良定义；自然映射 $\pi:A\to A/{\sim}$，$\pi(a)=[a]$，是满同态，核为 $\sim$。}{If $\sim$ is a congruence, then $[a]*[b]=[a*b]$ is independent of the choice of representatives, so the quotient algebra is well defined; the natural map $\pi:A\to A/{\sim}$, $\pi(a)=[a]$, is an epimorphism with kernel $\sim$.}
\PRF{若 $a\sim a'$、$b\sim b'$，由相容性 $a*b\sim a'*b'$，故 $[a*b]=[a'*b']$，运算良定义。$h$ 的核 $\ker h$ 恒为同余（$h(a)=h(b)$、$h(c)=h(d)$ 时 $h(a*c)=h(a)\oplus h(c)=h(b)\oplus h(d)=h(b*d)$），这是“同态—同余—商代数”三位一体的来源。}{If $a\sim a'$ and $b\sim b'$, compatibility gives $a*b\sim a'*b'$, so $[a*b]=[a'*b']$ and the operation is well defined. The kernel of any homomorphism is a congruence, since $h(a)=h(b)$ and $h(c)=h(d)$ imply $h(a*c)=h(a)\oplus h(c)=h(b)\oplus h(d)=h(b*d)$; this is the source of the homomorphism–congruence–quotient triad.}

\subsection{半群、独异点与群 | Semigroups, monoids and groups}
\DEF{只有结合律的代数叫\textbf{半群}；有单位元的半群叫\textbf{独异点}；每个元素都有逆元的独异点叫\textbf{群}。群 $(G,*)$ 满足：封闭、结合、存在单位元 $e$、每个 $a$ 存在逆元 $a^{-1}$；再满足交换律则称\textbf{交换群}（Abel 群）。群 $G$ 的元素个数 $|G|$ 叫\textbf{阶}；使 $a^k=e$ 的最小正整数 $k$ 叫元素 $a$ 的\textbf{阶}（不存在则称无限阶）。$G$ 的非空子集 $H$ 若对乘法与求逆封闭，称为\textbf{子群}，记 $H\le G$。$G=\langle g\rangle=\{g^k\mid k\in\mathbb Z\}$ 叫\textbf{循环群}，$g$ 是生成元。}{An algebra with associativity alone is a \textbf{semigroup}; a semigroup with an identity is a \textbf{monoid}; a monoid in which every element is invertible is a \textbf{group}. A group $(G,*)$ satisfies closure, associativity, an identity $e$ and an inverse $a^{-1}$ for every $a$; if it is also commutative it is \textbf{abelian}. The number $|G|$ is the \textbf{order} of the group, and the least positive $k$ with $a^k=e$ is the \textbf{order} of $a$ (infinite if none exists). A nonempty subset $H$ closed under the product and inverses is a \textbf{subgroup}, written $H\le G$. A \textbf{cyclic group} is $G=\langle g\rangle=\{g^k\mid k\in\mathbb Z\}$ with generator $g$.}
\THM{\textbf{子群判定}：$H\le G$ 当且仅当 $H\ne\varnothing$ 且 $\forall a,b\in H,\ ab^{-1}\in H$；有限情形下，只需 $H\ne\varnothing$ 且对乘法封闭。\textbf{循环群的性质}：循环群必交换；$n$ 阶循环群同构于 $(\mathbb Z_n,+_n)$，无限循环群同构于 $(\mathbb Z,+)$。}{$H\le G$ iff $H\ne\varnothing$ and $ab^{-1}\in H$ for all $a,b\in H$; in the finite case closure under the product alone suffices. A cyclic group is abelian; a cyclic group of order $n$ is isomorphic to $(\mathbb Z_n,+_n)$, and an infinite cyclic group to $(\mathbb Z,+)$.}
\PRF{子群判定：取 $a\in H$ 得 $e=aa^{-1}\in H$，再得 $a^{-1}=ea^{-1}\in H$，最后得乘法封闭。有限情形下，因 $a\in H$ 的幂 $a,a^2,\dots$ 有限，必重复，故 $a^k=e$ 对某 $k$，逆元 $a^{-1}=a^{k-1}\in H$。循环群交换：$g^mg^n=g^{m+n}=g^ng^m$。同构：映射 $g^k\mapsto k\bmod n$ 是 $n$ 阶循环群到 $\mathbb Z_n$ 的同构。}{Subgroup test: taking $a\in H$ gives $e=aa^{-1}\in H$, then $a^{-1}=ea^{-1}\in H$, then closure. In the finite case the powers $a,a^2,\dots$ repeat, so $a^k=e$ for some $k$ and $a^{-1}=a^{k-1}\in H$. Cyclic groups are abelian since $g^mg^n=g^{m+n}=g^ng^m$. Isomorphism: the map $g^k\mapsto k\bmod n$ is an isomorphism from a cyclic group of order $n$ to $\mathbb Z_n$.}

\subsection{陪集、Lagrange 定理与同态基本定理 | Cosets, Lagrange's theorem and the fundamental theorem of homomorphisms}
\DEF{设 $H\le G$。$aH=\{ah\mid h\in H\}$ 叫 $a$ 的\textbf{左陪集}，$Ha=\{ha\mid h\in H\}$ 叫\textbf{右陪集}；不同陪集的个数叫\textbf{指数} $[G:H]$。若 $\forall g\in G,\ gH=Hg$，则称 $H$ 为\textbf{正规子群}，记 $H\trianglelefteq G$。}{Let $H\le G$. The set $aH=\{ah\mid h\in H\}$ is a \textbf{left coset} and $Ha=\{ha\mid h\in H\}$ a \textbf{right coset}; the number of distinct cosets is the \textbf{index} $[G:H]$. If $gH=Hg$ for all $g\in G$, then $H$ is a \textbf{normal subgroup}, written $H\trianglelefteq G$.}
\THM{\textbf{Lagrange 定理}：设 $G$ 是有限群、$H\le G$，则 $|H|$ 整除 $|G|$，且 $|G|=|H|\cdot[G:H]$。推论：任何元素的阶整除 $|G|$，从而 $a^{|G|}=e$；素数阶群必为循环群（且只有平凡子群）。}{\textbf{Lagrange's theorem}: if $G$ is finite and $H\le G$, then $|H|$ divides $|G|$ and $|G|=|H|\cdot[G:H]$. Corollaries: the order of any element divides $|G|$, so $a^{|G|}=e$; a group of prime order is cyclic (and has only trivial subgroups).}
\PRF{左陪集构成 $G$ 的划分：$a\in aH$，故陪集覆盖 $G$；若 $aH\cap bH\ne\varnothing$，取 $x=ah_1=bh_2$，则对任意 $ah$ 有 $ah=bh_2h_1^{-1}h\in bH$，故 $aH\subseteq bH$，对称得 $bH\subseteq aH$，两陪集相同。又映射 $H\to aH$，$h\mapsto ah$，是双射，故 $|aH|=|H|$。于是 $|G|$ 等于陪集个数乘以 $|H|$。}{Left cosets partition $G$: $a\in aH$ so they cover $G$; if $aH\cap bH\ne\varnothing$, take $x=ah_1=bh_2$, then any $ah=bh_2h_1^{-1}h\in bH$, so $aH\subseteq bH$ and symmetrically $bH\subseteq aH$, hence the cosets are equal. The map $H\to aH$, $h\mapsto ah$, is a bijection, so $|aH|=|H|$. Thus $|G|$ equals the number of cosets times $|H|$.}
\THM{\textbf{群同态基本定理}：设 $\varphi:G\to G'$ 是群同态，则 $\ker\varphi\trianglelefteq G$，且诱导映射 $\overline\varphi:G/\ker\varphi\to\operatorname{Im}\varphi$，$\overline\varphi([a])=\varphi(a)$，是良定义的同构，即
\[
G/\ker\varphi\;\cong\;\operatorname{Im}\varphi .
\]
}{\textbf{Fundamental theorem of group homomorphisms}: for a group homomorphism $\varphi:G\to G'$, the kernel is normal in $G$ and the induced map $\overline\varphi:G/\ker\varphi\to\operatorname{Im}\varphi$ given by $\overline\varphi([a])=\varphi(a)$ is a well-defined isomorphism, i.e.\ $G/\ker\varphi\cong\operatorname{Im}\varphi$.}
\PRF{良定义：$[a]=[b]\Leftrightarrow ab^{-1}\in\ker\varphi$，故 $\varphi(a)=\varphi(b)$，$\overline\varphi$ 与代表元无关；它是满射到 $\operatorname{Im}\varphi$ 上，且 $\ker\overline\varphi=\{[e]\}$，故单射；保持运算显然，故为同构。核正规：$g(\ker\varphi)g^{-1}\subseteq\ker\varphi$，因 $\varphi(gkg^{-1})=\varphi(g)e\varphi(g)^{-1}=e$。}{Well-definedness: $[a]=[b]\Leftrightarrow ab^{-1}\in\ker\varphi$, so $\varphi(a)=\varphi(b)$ and $\overline\varphi$ is independent of representatives; it is surjective onto $\operatorname{Im}\varphi$ with trivial kernel $\{[e]\}$, hence injective; it preserves the operation, hence is an isomorphism. The kernel is normal since $\varphi(gkg^{-1})=\varphi(g)e\varphi(g)^{-1}=e$, so $g(\ker\varphi)g^{-1}\subseteq\ker\varphi$.}

\subsection{环与域 | Rings and fields}
\DEF{\textbf{环} $(R,+,\cdot)$ 满足：$(R,+)$ 是交换群，$(R,\cdot)$ 是半群，乘法对加法满足两条分配律。有乘法单位元的环叫\textbf{含幺环}，乘法交换的环叫\textbf{交换环}。含幺交换环且无零因子（$ab=0\Rightarrow a=0\ \text{或}\ b=0$）叫\textbf{整环}；每个非零元都有乘法逆元的含幺交换环叫\textbf{域}。}{A \textbf{ring} $(R,+,\cdot)$ satisfies: $(R,+)$ is an abelian group, $(R,\cdot)$ is a semigroup, and multiplication distributes over addition on both sides. A ring with a multiplicative identity is a \textbf{ring with unity}; a commutative one is a \textbf{commutative ring}. A commutative ring with unity and no zero divisors ($ab=0\Rightarrow a=0$ or $b=0$) is an \textbf{integral domain}; a commutative ring with unity in which every nonzero element is invertible is a \textbf{field}.}
\THM{\textbf{有限整环是域}：任何有限整环都是域。\textbf{$\mathbb Z_n$ 何时为域}：$\mathbb Z_n$ 是域当且仅当 $n$ 是素数。}{Every finite integral domain is a field. Moreover $\mathbb Z_n$ is a field exactly when $n$ is prime.}
\PRF{有限整环：取 $a\ne0$，映射 $x\mapsto ax$ 是单射（若 $ax=ay$ 则 $a(x-y)=0$，无零因子得 $x=y$），有限集上单射即满射，故存在 $b$ 使 $ab=1$。$\mathbb Z_n$：若 $n$ 合数，$n=rs$（$1<r,s<n$），则 $rs\equiv0$，$r,s$ 是零因子，非域；若 $n=p$ 素数，$\mathbb Z_p$ 是整环（$p\mid ab\Rightarrow p\mid a$ 或 $p\mid b$），且有限，故是域。}{Finite integral domain: for $a\ne0$ the map $x\mapsto ax$ is injective (if $ax=ay$ then $a(x-y)=0$, and no zero divisors force $x=y$); on a finite set injectivity is surjectivity, so some $b$ has $ab=1$. For $\mathbb Z_n$: if $n$ is composite, $n=rs$ with $1<r,s<n$, then $rs\equiv0$ and $r,s$ are zero divisors, so it is not a field; if $n=p$ is prime, $\mathbb Z_p$ is an integral domain ($p\mid ab\Rightarrow p\mid a$ or $p\mid b$) and finite, hence a field.}

\EX{在 $\mathbb Z_6$ 中，子群有 $\{0\}$、$\{0,3\}$、$\{0,2,4\}$、$\mathbb Z_6$，阶为 $1,2,3,6$，都整除 $6$，符合 Lagrange 定理。$H=\{0,3\}$ 的陪集为 $H$、$1+H=\{1,4\}$、$2+H=\{2,5\}$，指数 $[G:H]=3$，$|G|=|H|\cdot[G:H]=2\cdot3=6$。$\{0,3\}$ 是 $\mathbb Z_6$ 的（加法）子群，验证了子群判定：$3+3=0\in H$，故对加法封闭。}{In $\mathbb Z_6$ the subgroups are $\{0\}$, $\{0,3\}$, $\{0,2,4\}$ and $\mathbb Z_6$, of orders $1,2,3,6$, all dividing $6$, as Lagrange requires. The cosets of $H=\{0,3\}$ are $H$, $1+H=\{1,4\}$ and $2+H=\{2,5\}$, so the index $[G:H]=3$ and $|G|=|H|\cdot[G:H]=2\cdot3=6$. The set $\{0,3\}$ is an (additive) subgroup of $\mathbb Z_6$, verifying the subgroup test: $3+3=0\in H$, so it is closed under addition.}

\ASK{群同态基本定理说明了什么？}{What does the fundamental theorem of group homomorphisms say?}{同态的像同构于群模核的商群：$G/\ker\varphi\cong\operatorname{Im}\varphi$。直观上同态把核“压成”单位元，商掉核后剩下的信息与像完全一致；它把同态、同余、商代数统一起来，也解释了为什么核必须是正规子群。}{The image of a homomorphism is isomorphic to the quotient by its kernel: $G/\ker\varphi\cong\operatorname{Im}\varphi$. Intuitively the homomorphism collapses the kernel to the identity, and after quotienting by the kernel the remaining information matches the image exactly; it unifies homomorphisms, congruences and quotient algebras and explains why the kernel must be normal.}
\ASK{为什么 $\mathbb Z_n$ 是域当且仅当 $n$ 是素数？}{Why is $\mathbb Z_n$ a field exactly when $n$ is prime?}{$n$ 合数时有真因子 $r,s$ 使 $rs\equiv0$，出现零因子，非域；$n=p$ 素数时 $\mathbb Z_p$ 是有限整环（无零因子），而有限整环必为域，因为 $x\mapsto ax$ 单射即满射，每个非零元可逆。}{If $n$ is composite it has proper factors $r,s$ with $rs\equiv0$, giving zero divisors, so it is not a field; if $n=p$ is prime, $\mathbb Z_p$ is a finite integral domain (no zero divisors), and a finite integral domain is a field, since $x\mapsto ax$ is injective hence surjective and every nonzero element is invertible.}

\section[格与布尔代数（第 7 章）]{格与布尔代数（教材第 7 章）| Lattices and Boolean algebras (Chapter 7)}
\subsection{格的偏序定义 | The order-theoretic definition of a lattice}
\DEF{设 $(L,\preccurlyeq)$ 是偏序集。若任意 $a,b\in L$ 都有上确界 $a\vee b$ 与下确界 $a\wedge b$，则称 $L$ 为\textbf{格}，$\vee$ 叫\textbf{并}、$\wedge$ 叫\textbf{交}。若任意子集都有上确界与下确界，则称\textbf{完备格}。}{Let $(L,\preccurlyeq)$ be a poset. If every $a,b\in L$ have a least upper bound $a\vee b$ and a greatest lower bound $a\wedge b$, then $L$ is a \textbf{lattice}, with $\vee$ the \textbf{join} and $\wedge$ the \textbf{meet}. If every subset has a supremum and an infimum, it is a \textbf{complete lattice}.}
\THM{\textbf{格运算的基本律}：任意格满足幂等律、交换律、结合律、吸收律：
\[
a\vee a=a,\quad a\wedge a=a;\qquad a\vee b=b\vee a,\quad a\wedge b=b\wedge a;
\]
\[
(a\vee b)\vee c=a\vee(b\vee c),\quad (a\wedge b)\wedge c=a\wedge(b\wedge c);
\]
\[
a\vee(a\wedge b)=a,\quad a\wedge(a\vee b)=a,
\]
并且 $a\preccurlyeq b\Leftrightarrow a\vee b=b\Leftrightarrow a\wedge b=a$。}{Every lattice satisfies idempotence, commutativity, associativity and absorption, as displayed, together with $a\preccurlyeq b\Leftrightarrow a\vee b=b\Leftrightarrow a\wedge b=a$.}
\PRF{以吸收律为例：$a\wedge b\preccurlyeq a$，故 $a\vee(a\wedge b)$ 是 $a$ 与 $a\wedge b$ 的上确界；又 $a\wedge b\preccurlyeq a$，故该上确界就是 $a$。等价关系 $a\preccurlyeq b\Leftrightarrow a\vee b=b$：$b$ 是 $a,b$ 的上界，若 $a\preccurlyeq b$ 则 $b$ 是最小上界，反之 $a\preccurlyeq a\vee b=b$。}{For absorption: $a\wedge b\preccurlyeq a$, so $a\vee(a\wedge b)$ is the join of $a$ and $a\wedge b$; since $a\wedge b\preccurlyeq a$, that join is $a$ itself. For the equivalence $a\preccurlyeq b\Leftrightarrow a\vee b=b$: $b$ is an upper bound of $a,b$, and if $a\preccurlyeq b$ then $b$ is the least upper bound, while conversely $a\preccurlyeq a\vee b=b$.}
\REM{\textbf{对偶原理}：把偏序反向、$\vee$ 与 $\wedge$ 互换，格中任何真命题仍为真。并非每个偏序集都是格，例如四个元素 $a,b,c,d$ 满足 $a\prec c,a\prec d,b\prec c,b\prec d$ 时，$a,b$ 有上界 $c,d$ 但无最小上界。}{By the \textbf{duality principle}, reversing the order and swapping $\vee$ with $\wedge$ preserves every true statement. Not every poset is a lattice: with four elements $a,b,c,d$ and $a\prec c,a\prec d,b\prec c,b\prec d$, the pair $a,b$ has upper bounds $c,d$ but no least upper bound.}

\subsection{格的代数定义及其等价性 | The algebraic definition and its equivalence}
\DEF{若集合 $L$ 上的两个二元运算 $\vee,\wedge$ 满足交换律、结合律与吸收律，则称 $(L,\vee,\wedge)$ 为\textbf{代数格}。}{If two binary operations $\vee,\wedge$ on a set $L$ satisfy commutativity, associativity and absorption, then $(L,\vee,\wedge)$ is an \textbf{algebraic lattice}.}
\THM{\textbf{两种定义的等价性}：偏序格与代数格是同一个概念。具体地：(i) 偏序格 $(L,\preccurlyeq)$ 的并、交满足交换、结合、吸收，故是代数格；(ii) 代数格上定义 $a\preccurlyeq b\Leftrightarrow a\wedge b=a$，则 $\preccurlyeq$ 是偏序，且 $a\wedge b$、$a\vee b$ 分别是该偏序下的下确界与上确界。}{The two definitions coincide. Precisely: (i) the join and meet of an order-theoretic lattice satisfy commutativity, associativity and absorption, so it is an algebraic lattice; (ii) on an algebraic lattice, defining $a\preccurlyeq b\Leftrightarrow a\wedge b=a$ makes $\preccurlyeq$ a partial order in which $a\wedge b$ and $a\vee b$ are the infimum and supremum.}
\PRF{(i) 已由格运算基本律给出。(ii) 自反性由幂等律 $a\wedge a=a$；反对称性由交换律：$a\wedge b=a$ 且 $b\wedge a=b$ 得 $a=b$；传递性由结合律：$a\wedge b=a$、$b\wedge c=b$ 得 $a\wedge c=(a\wedge b)\wedge c=a\wedge(b\wedge c)=a\wedge b=a$。下确界：$a\wedge b$ 是 $a,b$ 的下界（吸收律给出 $(a\wedge b)\wedge a=a\wedge b$），且任何下界 $d$ 满足 $d=d\wedge(a\wedge b)\preccurlyeq a\wedge b$；上确界对偶。}{(i) is the basic laws of lattice operations. (ii) Reflexivity from $a\wedge a=a$; antisymmetry from commutativity: $a\wedge b=a$ and $b\wedge a=b$ give $a=b$; transitivity from associativity: $a\wedge b=a$ and $b\wedge c=b$ give $a\wedge c=(a\wedge b)\wedge c=a\wedge(b\wedge c)=a\wedge b=a$. For the infimum, $a\wedge b$ is a lower bound of $a,b$ (absorption gives $(a\wedge b)\wedge a=a\wedge b$), and any lower bound $d$ satisfies $d=d\wedge(a\wedge b)\preccurlyeq a\wedge b$; the supremum is dual.}

\subsection{分配格、模格与有补格 | Distributive, modular and complemented lattices}
\DEF{格 $L$ 若满足分配律 $a\wedge(b\vee c)=(a\wedge b)\vee(a\wedge c)$（等价于其对偶式），称为\textbf{分配格}。满足模律“$a\preccurlyeq b\Rightarrow a\vee(c\wedge b)=(a\vee c)\wedge b$”的格叫\textbf{模格}。有最小元 $0$ 与最大元 $1$ 的格叫\textbf{有界格}；有界格中，若 $a\wedge b=0$ 且 $a\vee b=1$，则 $b$ 叫 $a$ 的\textbf{补元}；每个元素都有补元的格叫\textbf{有补格}。}{A lattice satisfying the distributive law $a\wedge(b\vee c)=(a\wedge b)\vee(a\wedge c)$ (equivalently its dual) is \textbf{distributive}. One satisfying the modular law ``$a\preccurlyeq b\Rightarrow a\vee(c\wedge b)=(a\vee c)\wedge b$'' is \textbf{modular}. A lattice with least element $0$ and greatest element $1$ is \textbf{bounded}; in a bounded lattice, $b$ is a \textbf{complement} of $a$ if $a\wedge b=0$ and $a\vee b=1$; a lattice in which every element has a complement is \textbf{complemented}.}
\THM{\textbf{分配与模}：分配格必是模格，反之不然。两个经典非分配格是\textbf{菱形} $M_3$（三个原子两两不可比，共用一个底与一个顶）与\textbf{五边形} $N_5$；$M_3$ 是模格但非分配格，$N_5$ 连模律都不满足。\textbf{Birkhoff 定理}：格是分配格当且仅当它不含与 $M_3$ 或 $N_5$ 同构的子格。}{Every distributive lattice is modular, but not conversely. The two classical non-distributive lattices are the \textbf{diamond} $M_3$ (three pairwise incomparable atoms sharing one bottom and one top) and the \textbf{pentagon} $N_5$; $M_3$ is modular but not distributive, and $N_5$ fails even the modular law. \textbf{Birkhoff's theorem}: a lattice is distributive iff it has no sublattice isomorphic to $M_3$ or $N_5$.}
\PRF{分配 $\Rightarrow$ 模：设 $a\preccurlyeq b$，则 $a\wedge b=a$，于是 $(a\vee c)\wedge b=(a\wedge b)\vee(c\wedge b)=a\vee(c\wedge b)$。$M_3$ 中三个原子 $p,q,r$ 满足 $p\wedge(q\vee r)=p\wedge1=p$，而 $(p\wedge q)\vee(p\wedge r)=0\vee0=0$，故不分配；但它满足模律（可逐一验证）。}{Distributive $\Rightarrow$ modular: if $a\preccurlyeq b$ then $a\wedge b=a$, so $(a\vee c)\wedge b=(a\wedge b)\vee(c\wedge b)=a\vee(c\wedge b)$. In $M_3$ the three atoms $p,q,r$ give $p\wedge(q\vee r)=p\wedge1=p$ while $(p\wedge q)\vee(p\wedge r)=0\vee0=0$, so it is not distributive, yet it satisfies the modular law (checkable case by case).}
\THM{\textbf{分配格中补元唯一}：在有界分配格中，元素若存在补元则唯一。}{In a bounded distributive lattice, a complement of an element, if it exists, is unique.}
\PRF{设 $b,b'$ 都是 $a$ 的补元，则
\[
b=b\wedge1=b\wedge(a\vee b')=(b\wedge a)\vee(b\wedge b')=0\vee(b\wedge b')=b\wedge b',
\]
对称得 $b'=b\wedge b'$，故 $b=b'$。注意一般有补格中补元不必唯一（$M_3$ 中每个原子有两个补元），唯一性依赖分配律。}{If $b,b'$ both complement $a$, then the displayed chain gives $b=b\wedge b'$, and symmetrically $b'=b\wedge b'$, hence $b=b'$. In a general complemented lattice complements need not be unique (in $M_3$ each atom has two complements); uniqueness relies on distributivity.}

\subsection{布尔代数 | Boolean algebras}
\DEF{既有补又分配的格叫\textbf{布尔代数}，记 $(B,\vee,\wedge,{}',0,1)$，$'$ 为补运算。等价地可用 \textbf{Huntington 公理}：$\vee,\wedge$ 交换且分配；存在 $0,1$ 使 $a\vee0=a$、$a\wedge1=a$；且 $a\vee a'=1$、$a\wedge a'=0$。}{A lattice that is both complemented and distributive is a \textbf{Boolean algebra}, written $(B,\vee,\wedge,{}',0,1)$ with $'$ the complement. Equivalently, the \textbf{Huntington axioms} hold: $\vee,\wedge$ are commutative and distributive; there are $0,1$ with $a\vee0=a$ and $a\wedge1=a$; and $a\vee a'=1$, $a\wedge a'=0$.}
\THM{\textbf{布尔代数的基本恒等式}：对一切 $a,b$ 成立德摩根律 $(a\vee b)'=a'\wedge b'$、$(a\wedge b)'=a'\vee b'$，以及 $0'=1$、$1'=0$、$(a')'=a$、$a\vee(a'\wedge b)=a\vee b$。}{For all $a,b$ in a Boolean algebra, De Morgan's laws hold, $(a\vee b)'=a'\wedge b'$ and $(a\wedge b)'=a'\vee b'$, together with $0'=1$, $1'=0$, $(a')'=a$ and $a\vee(a'\wedge b)=a\vee b$.}
\THM{\textbf{Stone 表示定理}：每个布尔代数都同构于某个集合代数（集合域，即某集合的幂集在 $\cup,\cap,{}^c$ 下构成的布尔代数）；有限布尔代数同构于 $(\mathcal P(A),\cup,\cap,{}^c)$，且 $|B|=2^{|A|}$（其中 $A$ 是原子集）。}{Every Boolean algebra is isomorphic to an algebra of sets (a field of sets, i.e.\ a Boolean algebra of subsets under $\cup,\cap,{}^c$); a finite Boolean algebra is isomorphic to $(\mathcal P(A),\cup,\cap,{}^c)$ with $|B|=2^{|A|}$, where $A$ is the set of atoms.}
\PRF{有限情形的思路：取 $B$ 的原子（盖住 $0$ 的元素）为 $A$，则每个元素 $x$ 唯一地写成若干原子之并，映射 $x\mapsto\{a\in A\mid a\preccurlyeq x\}$ 给出到 $\mathcal P(A)$ 的同构；因此有限布尔代数的阶必为 $2$ 的幂。}{The idea in the finite case: take the atoms (elements covering $0$) as $A$; every element $x$ is uniquely a join of atoms, and the map $x\mapsto\{a\in A\mid a\preccurlyeq x\}$ is an isomorphism to $\mathcal P(A)$; hence the order of a finite Boolean algebra is a power of $2$.}
\REM{\textbf{布尔函数}由变元与 $0,1,\vee,\wedge,{}'$ 构成；由第 1 章的主范式定理，每个布尔函数都有唯一的主析取范式，故布尔代数与命题逻辑是同一对象——命题变元即开关变量，$\vee,\wedge,{}'$ 即并联、串联、反相。逻辑电路、数据库查询与集合运算都可用它统一描述。}{A \textbf{Boolean function} is built from variables, $0,1,\vee,\wedge,{}'$; by the principal normal form theorem of Chapter 1, every Boolean function has a unique PDNF, so Boolean algebras and propositional logic are the same object: propositional variables are switch variables and $\vee,\wedge,{}'$ are parallel, series and inversion. Logic circuits, database queries and set operations are all described by it.}

\EX{取 $n=30=2\cdot3\cdot5$ 的正因子集 $D=\{1,2,3,5,6,10,15,30\}$，在 $\gcd$（交）与 $\operatorname{lcm}$（并）下构成布尔代数，$0=1$、$1=30$，补元为 $30/d$。例如 $6$ 的补元是 $5$，因为 $\gcd(6,5)=1$、$\operatorname{lcm}(6,5)=30$。原子是 $2,3,5$，故 $D$ 同构于 $\mathcal P(\{2,3,5\})$，$|D|=2^3=8$，符合 Stone 表示定理。}{The divisor set $D=\{1,2,3,5,6,10,15,30\}$ of $n=30=2\cdot3\cdot5$ under $\gcd$ (meet) and $\operatorname{lcm}$ (join) is a Boolean algebra with $0=1$, $1=30$ and complement $30/d$. The complement of $6$ is $5$, since $\gcd(6,5)=1$ and $\operatorname{lcm}(6,5)=30$. The atoms are $2,3,5$, so $D\cong\mathcal P(\{2,3,5\})$ with $|D|=2^3=8$, as Stone's representation predicts.}

\ASK{格的两个定义为什么等价？}{Why are the two definitions of a lattice equivalent?}{偏序定义给出任意二元的上、下确界，直接验证满足交换、结合、吸收；反过来在代数格上以 $a\preccurlyeq b\Leftrightarrow a\wedge b=a$ 定义序，吸收律保证它是偏序，等式保证 $a\wedge b$、$a\vee b$ 分别是下、上确界。于是“序的语言”与“运算的语言”可互换。}{The order-theoretic definition gives binary suprema and infima, which one verifies satisfy commutativity, associativity and absorption; conversely, on an algebraic lattice define $a\preccurlyeq b\Leftrightarrow a\wedge b=a$, use absorption to show it is a partial order, and the identities to show $a\wedge b$ and $a\vee b$ are the infimum and supremum. The languages of order and of operations are interchangeable.}
\ASK{为什么布尔代数要求“有补”与“分配”两个条件？}{Why does a Boolean algebra require both complemented and distributive?}{有补保证每个元素有“否定”，分配保证补元唯一并使德摩根律成立；只有有补（如 $M_3$）时补元不唯一，无法定义确定的补运算；只有分配而不有补（如链）时缺少否定。两者合起来才与命题逻辑、集合代数、开关代数同构。}{Complementedness gives a ``negation'' for every element, and distributivity makes complements unique and yields De Morgan's laws; with complementedness alone (as in $M_3$) complements are not unique and no well-defined complement operation exists, while with distributivity alone (as in a chain) negation is missing. Together they match propositional logic, set algebra and switching algebra.}

\section[图论（第 8 章）]{图论（教材第 8 章）| Graph theory (Chapter 8)}
\subsection{图的基本概念 | Basic notions of graphs}
\DEF{\textbf{无向图} $G=(V,E)$ 由顶点集 $V$ 与边集 $E$（$V$ 的无序对）组成；\textbf{有向图}的边是有序对。无重边无自回路的图叫\textbf{简单图}。顶点 $v$ 的\textbf{度} $\deg v$ 是与其关联的边数（自回路计 $2$）；度为 $0$ 的顶点叫孤立点。两图\textbf{同构}当且仅当存在保持邻接关系的双射 $\varphi:V_1\to V_2$。}{An \textbf{undirected graph} $G=(V,E)$ consists of a vertex set $V$ and an edge set $E$ of unordered pairs; in a \textbf{directed graph} the edges are ordered pairs. A \textbf{simple graph} has no multiple edges and no loops. The \textbf{degree} $\deg v$ of a vertex $v$ is the number of incident edges (a loop counting $2$); a vertex of degree $0$ is isolated. Two graphs are \textbf{isomorphic} iff there is a bijection $\varphi:V_1\to V_2$ preserving adjacency.}
\THM{\textbf{握手定理}：对任意图 $G=(V,E)$，$\sum_{v\in V}\deg v=2|E|$；有向图中 $\sum_v d^+(v)=\sum_v d^-(v)=|E|$。推论：奇度顶点的个数必为偶数。}{For any graph $G=(V,E)$, $\sum_{v\in V}\deg v=2|E|$; in a digraph $\sum_v d^+(v)=\sum_v d^-(v)=|E|$. Corollary: the number of odd-degree vertices is even.}
\PRF{每条边有两个端点，求和时恰被计数两次；对自回路也贡献 $2$。故度数和等于 $2|E|$，是偶数，奇度顶点个数只能为偶数。}{Each edge has two endpoints, so summing degrees counts every edge exactly twice (a loop also contributes $2$). Hence the degree sum is $2|E|$, an even number, so the number of odd-degree vertices is even.}

\subsection{路径、回路与连通性 | Paths, circuits and connectivity}
\DEF{顶点与边交替的序列叫\textbf{途径}；边不重复的叫\textbf{迹}；顶点不重复的叫\textbf{路径}；起点等于终点的路径叫\textbf{回路}。无向图任意两点有路径相连则称\textbf{连通}，极大连通子图叫\textbf{连通分支}；删去后使分支数增加的顶点叫\textbf{割点}、边叫\textbf{桥}。经过每条边恰一次的迹叫 \textbf{Euler 迹}，闭的 Euler 迹叫 \textbf{Euler 回路}；经过每个顶点恰一次的路径叫 \textbf{Hamilton 路径}，闭的 Hamilton 路径叫 \textbf{Hamilton 回路}。}{An alternating sequence of vertices and edges is a \textbf{walk}; one with no repeated edge a \textbf{trail}; one with no repeated vertex a \textbf{path}; a path whose endpoints coincide a \textbf{circuit}. An undirected graph is \textbf{connected} when any two vertices are joined by a path; maximal connected subgraphs are \textbf{components}; a vertex (edge) whose deletion increases the number of components is a \textbf{cut vertex} (\textbf{bridge}). A trail using every edge exactly once is an \textbf{Euler trail}, a closed one an \textbf{Euler circuit}; a path visiting every vertex exactly once is a \textbf{Hamiltonian path}, a closed one a \textbf{Hamiltonian cycle}.}
\THM{\textbf{Euler 定理}：连通无向图存在 Euler 回路当且仅当所有顶点度为偶数；存在 Euler 路径（非回路）当且仅当恰有两个奇度顶点（它们为端点）。连通有向图存在 Euler 回路当且仅当每个顶点入度等于出度。}{A connected undirected graph has an Euler circuit iff every degree is even; it has an Euler path (but no circuit) iff exactly two vertices have odd degree (the endpoints). A connected digraph has an Euler circuit iff the in-degree equals the out-degree at every vertex.}
\PRF{必要性显然：Euler 回路每进入一个顶点必离开一次，故每个顶点度为偶数。充分性：从任一顶点出发沿未用边走，直到不能再走，所停顶点必为出发点（否则奇度）；把得到的回路“拼接”进尚未用尽边的顶点的子回路，反复直到覆盖所有边。Hamilton 问题没有类似的简单充要条件，是 NP 完全问题。}{Necessity is clear: an Euler circuit enters and leaves each vertex equally often, so every degree is even. Sufficiency: start from any vertex and follow unused edges until forced to stop; the stop vertex must be the start (otherwise odd degree); splice the subcircuits at vertices with unused edges, repeating until all edges are used. The Hamiltonian problem has no such simple necessary-and-sufficient condition and is NP-complete.}

\subsection{图的矩阵表示 | Matrix representations of graphs}
\DEF{无向简单图 $G$ 的\textbf{邻接矩阵} $A=(a_{ij})$ 满足 $a_{ij}=1$ 当且仅当 $v_i$ 与 $v_j$ 相邻（有向图 $a_{ij}=1$ 表 $v_i\to v_j$）；\textbf{关联矩阵} $M=(m_{ij})$ 记录顶点与边的关联（无向图 $m_{ij}=1$ 当且仅当 $v_i$ 是 $e_j$ 的端点）。}{The \textbf{adjacency matrix} $A=(a_{ij})$ of an undirected simple graph has $a_{ij}=1$ iff $v_i$ and $v_j$ are adjacent (in a digraph $a_{ij}=1$ means $v_i\to v_j$); the \textbf{incidence matrix} $M=(m_{ij})$ records vertex-edge incidence ($m_{ij}=1$ iff $v_i$ is an endpoint of $e_j$ in an undirected graph).}
\THM{\textbf{邻接矩阵的幂}：$(A^k)_{ij}$ 等于从 $v_i$ 到 $v_j$ 长度为 $k$ 的途径条数。把布尔运算换成 $\vee,\wedge$，则\textbf{可达矩阵} $P=A^{(1)}\vee\cdots\vee A^{(n)}$ 的 $(i,j)$ 元为 $1$ 当且仅当 $v_i$ 可达 $v_j$。}{$(A^k)_{ij}$ equals the number of walks of length $k$ from $v_i$ to $v_j$. With Boolean operations ($\vee,\wedge$), the \textbf{reachability matrix} $P=A^{(1)}\vee\cdots\vee A^{(n)}$ has $(i,j)$-entry $1$ iff $v_i$ reaches $v_j$.}
\PRF{对 $k$ 归纳：$(A^k)_{ij}=\sum_\ell(A^{k-1})_{i\ell}\,a_{\ell j}$，即长度 $k$ 的途径按倒数第二个顶点分类，是长度 $k-1$ 的途径接上一条边；这正是矩阵乘法的定义。}{Induct on $k$: $(A^k)_{ij}=\sum_\ell(A^{k-1})_{i\ell}\,a_{\ell j}$, classifying walks of length $k$ by the penultimate vertex, a walk of length $k-1$ followed by an edge; this is exactly the definition of matrix multiplication.}

\subsection{二部图与匹配 | Bipartite graphs and matching}
\DEF{顶点集可分拆为 $X,Y$，使每条边一端在 $X$、另一端在 $Y$，则称 $G$ 为\textbf{二部图}；$X,Y$ 之间每对顶点都相邻的二部图叫\textbf{完全二部图} $K_{m,n}$。边集 $M$ 若两两不共端点，称为\textbf{匹配}；饱和 $X$（每个 $x\in X$ 都是某匹配边的端点）的匹配叫\textbf{饱和 $X$ 的匹配}；饱和所有顶点的匹配叫\textbf{完美匹配}。}{A graph is \textbf{bipartite} when its vertices split into $X,Y$ so that every edge has one endpoint in $X$ and one in $Y$; if every pair across $X,Y$ is adjacent it is the \textbf{complete bipartite graph} $K_{m,n}$. A set $M$ of pairwise disjoint edges is a \textbf{matching}; one \textbf{saturating $X$} makes every $x\in X$ an endpoint; one saturating every vertex is a \textbf{perfect matching}.}
\THM{\textbf{二部图判定}：无向图是二部图当且仅当它不含奇长度回路。}{An undirected graph is bipartite iff it contains no cycle of odd length.}
\PRF{必要性：二部图每条边都跨 $X,Y$，故沿任何回路必在 $X,Y$ 间交替，长度只能为偶数。充分性：对每个连通分支任取一点放入 $X$，用 BFS 按奇偶层交替染 $X,Y$；无奇回路保证染色不冲突，得二部划分。}{Necessity: in a bipartite graph every edge crosses $X,Y$, so any cycle alternates between them and has even length. Sufficiency: in each component put an arbitrary vertex in $X$ and colour by BFS parity; the absence of odd cycles guarantees a consistent colouring, giving the bipartition.}
\THM{\textbf{Hall 定理（婚姻定理）}：设 $G=(X\cup Y,E)$ 是二部图，$N(S)$ 表 $S\subseteq X$ 的邻点集。存在饱和 $X$ 的匹配当且仅当对每个 $S\subseteq X$ 有 $|N(S)|\ge|S|$。}{Let $G=(X\cup Y,E)$ be bipartite and $N(S)$ the set of neighbours of $S\subseteq X$. There is a matching saturating $X$ iff $|N(S)|\ge|S|$ for every $S\subseteq X$.}
\PRF{必要性显然：$S$ 中 $|S|$ 个点要匹配到不同的邻点。充分性用\textbf{增广路径}：取未饱和 $X$ 的最大匹配 $M$，从某未饱和点 $x_0$ 出发沿“未匹配边到 $Y$、匹配边回 $X$”交错前进，得交错树；Hall 条件保证树中 $Y$ 侧顶点严格多于与 $x_0$ 同侧的 $X$ 顶点，故必有一未饱和的 $y$ 可达，沿该交错路径翻转（匹配边与非匹配边互换）使匹配大小 $+1$，与最大性矛盾。}{Necessity is clear: $|S|$ vertices need distinct neighbours. Sufficiency uses \textbf{augmenting paths}: take a maximum matching $M$ not saturating $X$, start from an unsaturated $x_0$ and grow an alternating tree (unmatched edges into $Y$, matched edges back into $X$); Hall's condition forces strictly more $Y$-vertices than $X$-vertices on the $x_0$ side, so some unsaturated $y$ is reachable, and flipping along that alternating path increases the matching by $1$, contradicting maximality.}

\subsection{平面图与图的着色 | Planar graphs and colouring}
\DEF{图能画在平面上使任意两条边只在公共端点相交，则称\textbf{平面图}；这样的画法把平面分成若干\textbf{面}，无界面叫外部面。顶点着色是给顶点分配颜色使相邻顶点异色；最少颜色数叫\textbf{色数} $\chi(G)$。}{A graph is \textbf{planar} if it can be drawn in the plane so that edges meet only at common endpoints; such a drawing divides the plane into \textbf{faces}, the unbounded one being the outer face. A \textbf{vertex colouring} assigns colours so that adjacent vertices differ; the least number of colours is the \textbf{chromatic number} $\chi(G)$.}
\THM{\textbf{Euler 公式}：连通平面图满足 $V-E+F=2$，其中 $V,E,F$ 是顶点数、边数、面数（含外部面）。推论：若 $G$ 是 $V\ge3$ 的简单平面图，则 $E\le3V-6$；若 $G$ 还是二部平面图，则 $E\le2V-4$。因此 $K_5$ 与 $K_{3,3}$ 都不是平面图。}{For a connected planar graph, $V-E+F=2$, where $V,E,F$ are the numbers of vertices, edges and faces (including the outer face). Corollary: a simple planar graph with $V\ge3$ satisfies $E\le3V-6$; a bipartite planar graph satisfies $E\le2V-4$. Hence $K_5$ and $K_{3,3}$ are nonplanar.}
\PRF{Euler 公式对 $E$ 归纳：取一棵生成树（$V-1$ 条边，$F=1$，公式成立），再逐条加回树外之边，每加一条边恰好增加一个面，$E,F$ 同步 $+1$，公式保持。$E\le3V-6$：每个面至少由 $3$ 条边围成，故 $3F\le2E$（每条边至多属于两个面），代入 $F=2-V+E$ 得 $E\le3V-6$；二部图无三角面，每个面至少 $4$ 条边，得 $2E\ge4F$，即 $E\le2V-4$。$K_5$ 有 $V=5,E=10>9$；$K_{3,3}$ 有 $V=6,E=9>8$，均矛盾。}{Euler's formula is proved by induction on $E$: take a spanning tree ($V-1$ edges, $F=1$, the formula holds) and add back the non-tree edges one by one; each added edge increases both $E$ and $F$ by $1$, preserving the formula. For $E\le3V-6$: each face is bounded by at least $3$ edges, so $3F\le2E$ (each edge borders at most two faces); substituting $F=2-V+E$ gives $E\le3V-6$. A bipartite planar graph has no triangular face, so each face has at least $4$ edges, giving $2E\ge4F$, i.e.\ $E\le2V-4$. $K_5$ has $V=5,E=10>9$, and $K_{3,3}$ has $V=6,E=9>8$, both contradictions.}
\THM{\textbf{Kuratowski 定理（陈述）}：图是平面图当且仅当它不含与 $K_5$ 或 $K_{3,3}$ 同胚的子图。\textbf{四色定理（陈述）}：任何平面图的色数 $\chi(G)\le4$。}{A graph is planar iff it contains no subgraph homeomorphic to $K_5$ or $K_{3,3}$. The \textbf{four colour theorem}: every planar graph has $\chi(G)\le4$.}
\REM{基本色数：$\chi(K_n)=n$；$\chi(C_n)=2$（$n$ 偶）或 $3$（$n$ 奇）；图是二部图当且仅当 $\chi\le2$。四色定理的已知证明依赖计算机穷举约 $1900$ 种不可约构形；着色建模地图着色、排课、频率分配等。}{Basic chromatic numbers: $\chi(K_n)=n$; $\chi(C_n)=2$ for even $n$ and $3$ for odd $n$; a graph is bipartite iff $\chi\le2$. The known proof of the four colour theorem relies on computer enumeration of about $1900$ reducible configurations; colouring models map colouring, timetabling and frequency assignment.}

\subsection{树与最小生成树 | Trees and minimum spanning trees}
\DEF{\textbf{树}是连通且无回路的无向图；树的度为 $1$ 的顶点叫\textbf{叶}。连通图 $G$ 的\textbf{生成树}是包含全部顶点、边集为 $E(G)$ 子集的树。带权图的\textbf{最小生成树}（MST）是总权最小的生成树。}{A \textbf{tree} is a connected acyclic undirected graph; its degree-$1$ vertices are \textbf{leaves}. A \textbf{spanning tree} of a connected graph $G$ is a tree containing all vertices whose edge set is a subset of $E(G)$. A \textbf{minimum spanning tree} (MST) of a weighted graph is a spanning tree of least total weight.}
\THM{\textbf{树的等价刻画}：设 $G$ 是 $n$ 个顶点的图，则下列命题等价：(1) $G$ 是树（连通且无回路）；(2) $G$ 连通且恰有 $n-1$ 条边；(3) $G$ 无回路且恰有 $n-1$ 条边；(4) $G$ 任意两顶点间恰有一条路径；(5) $G$ 连通且每条边都是桥；(6) $G$ 无回路，但添加任意一条新边就产生唯一回路。}{Let $G$ be a graph on $n$ vertices. The following are equivalent: (1) $G$ is a tree (connected and acyclic); (2) $G$ is connected with exactly $n-1$ edges; (3) $G$ is acyclic with exactly $n-1$ edges; (4) between any two vertices there is exactly one path; (5) $G$ is connected and every edge is a bridge; (6) $G$ is acyclic, but adding any new edge creates a unique cycle.}
\PRF{思路是沿环 (1)$\Rightarrow$(2)$\Rightarrow$(3)$\Rightarrow$(4)$\Rightarrow$(5)$\Rightarrow$(6)$\Rightarrow$(1) 推：(1)$\Rightarrow$(2) 用归纳（每棵树至少两片叶，删叶使 $n$ 与 $E$ 同时减 $1$）；(2)$\Rightarrow$(3) 若含回路，删去回路上一条边仍连通，重复到无回路得一棵树，其边数 $<n-1$，矛盾；(3)$\Rightarrow$(4) 若两点有两条路径则含回路，若有孤立点则边数不足；(4)$\Rightarrow$(5)(6) 与 (6)$\Rightarrow$(1) 直接验证。}{The proof follows the cycle (1)$\Rightarrow$(2)$\Rightarrow$(3)$\Rightarrow$(4)$\Rightarrow$(5)$\Rightarrow$(6)$\Rightarrow$(1): (1)$\Rightarrow$(2) by induction (a tree has at least two leaves, and deleting a leaf reduces $n$ and $E$ together); (2)$\Rightarrow$(3): if a cycle existed, delete an edge of it and stay connected, repeat to get a tree, which has fewer than $n-1$ edges, a contradiction; (3)$\Rightarrow$(4): two paths between two vertices form a cycle, and an isolated vertex would make the edge count too small; (4)$\Rightarrow$(5), (6) and (6)$\Rightarrow$(1) are direct checks.}
\THM{\textbf{生成树的存在与割性质}：每个连通图都有生成树。设 $S$ 是带权图中把顶点分成两部分的\textbf{割}，则跨越 $S$ 的权最小边必属于某棵最小生成树。}{Every connected graph has a spanning tree. Let a \textbf{cut} split the vertices of a weighted graph into two parts; a least-weight edge crossing the cut belongs to some minimum spanning tree.}
\PRF{生成树：反复删除回路上的一条边直到无回路。割性质：若最小边 $e$ 不在某 MST $T$ 中，把它加入 $T$ 得唯一回路，该回路必含另一条跨割边 $e'$；由 $e$ 最小，$w(e)\le w(e')$，用 $e$ 替换 $e'$ 得到权不减（实为不增）的新生成树，故存在含 $e$ 的 MST。\textbf{Kruskal} 按权从小到大加入不构成回路的边，\textbf{Prim} 从一点出发每次取连接已选与未选集合的最小边，两者都是割性质的直接实现。}{Spanning tree: repeatedly delete an edge on a cycle until acyclic. Cut property: if a least edge $e$ is not in some MST $T$, adding it to $T$ forms a unique cycle, which contains another edge $e'$ crossing the cut; since $e$ is least, $w(e)\le w(e')$, and replacing $e'$ by $e$ gives a new spanning tree of weight no larger, so some MST contains $e$. \textbf{Kruskal} adds edges in increasing weight unless they form a cycle, and \textbf{Prim} grows from one vertex by always taking the least edge joining the chosen to the unchosen set; both implement the cut property.}

\subsection{有向树与根树 | Directed trees and rooted trees}
\DEF{有向图中，若存在一个顶点（\textbf{根}）入度为 $0$、其余顶点入度均为 $1$，且从根到每个顶点都有有向路径，则称其为\textbf{有向树}（树形图，arborescence）。根树中从根到顶点的路径长度叫\textbf{层数}，最大层数叫\textbf{高度}；顶点的直接后继叫儿子、前驱叫父亲。每个顶点至多 $m$ 个儿子的根树叫 $m$ 叉树，$m=2$ 时为二叉树。}{In a digraph, if one vertex (the \textbf{root}) has in-degree $0$, every other vertex has in-degree $1$, and there is a directed path from the root to every vertex, the digraph is a \textbf{directed tree} (arborescence). In a rooted tree, the length of the path from the root to a vertex is its \textbf{level}, and the maximum level is the \textbf{height}; the immediate successors of a vertex are its children and its predecessor is its parent. A rooted tree with at most $m$ children per vertex is an $m$-ary tree, a binary tree when $m=2$.}

\EX{对六个顶点 $a,b,c,d,e,f$ 的带权图，Kruskal 按权升序扫描并舍去成回路的边：取 $ab(1),bc(2)$，$ac(2)$ 与 $a,b,c$ 成回路舍去，取 $cd(3),de(4),ef(5)$，此时已有 $5=n-1$ 条边且连通全部顶点，停止。MST 为 $\{ab,bc,cd,de,ef\}$，总权 $1+2+3+4+5=15$。}{For a weighted graph on six vertices $a,b,c,d,e,f$, Kruskal scans edges in increasing weight and rejects those forming a cycle: take $ab(1),bc(2)$, reject $ac(2)$ (it closes $a{-}b{-}c$), take $cd(3),de(4),ef(5)$; now there are $5=n-1$ edges connecting all vertices, so it stops. The MST is $\{ab,bc,cd,de,ef\}$ with total weight $1+2+3+4+5=15$.}

\ASK{握手定理为什么成立？能推出什么？}{Why does the handshake theorem hold, and what follows?}{每条边有两个端点，求和时每条边恰被计数两次，故 $\sum\deg v=2|E|$。于是度数和必为偶数，奇度顶点个数必为偶数；由此可快速排除“5 个顶点全为 3 度”这类不可能的度序列。}{Every edge has two endpoints, so summing degrees counts each edge twice, hence $\sum\deg v=2|E|$. The degree sum is therefore even and the number of odd-degree vertices is even; this quickly rules out impossible degree sequences such as five vertices all of degree $3$.}
\ASK{Euler 公式需要什么条件？为什么 $K_5$ 非平面？}{What are the hypotheses of Euler's formula, and why is $K_5$ nonplanar?}{Euler 公式 $V-E+F=2$ 要求图连通且平面，$F$ 含外部面；不连通时为 $V-E+F=1+c$（$c$ 为分支数）。简单平面图每个面至少三条边，故 $3F\le2E$，代入得 $E\le3V-6$；$K_5$ 有 $V=5,E=10>9$，矛盾，故非平面。}{Euler's formula $V-E+F=2$ requires a connected planar graph with $F$ counting the outer face; for a disconnected graph it becomes $V-E+F=1+c$ ($c$ components). In a simple planar graph each face has at least three edges, so $3F\le2E$ and substitution gives $E\le3V-6$; $K_5$ has $V=5,E=10>9$, a contradiction, so it is nonplanar.}

\section*{实际应用与面试表述 | Applications and interview presentation}
\APS{本讲义刻意把应用压成要点：命题逻辑与布尔代数统一为开关电路与数字逻辑（与或非门实现联结词，主析取范式给“积之和”实现，$\{\uparrow\}$ 的完备性说明 NAND 门即可搭建一切组合电路，SAT 求解器用于硬件验证与规划）；关系是关系数据库与依赖/分组的模型（传递闭包回答“可达”“上级的上级”）；群环域支撑编码与密码（$\mathbb Z_n$ 是 RSA 与离散对数的舞台，线性码是 $\mathbb F_2^n$ 的子空间）；图论支撑最短路、最小生成树、最大流、二分匹配与着色。面试表述建议分三层：语言层（逻辑与集合）、结构层（关系、偏序、等价类）、算法层（闭包、匹配、生成树、着色、拓扑排序），每层各举一个具体名词即可。}{This review deliberately compresses applications into essentials. Propositional logic and Boolean algebra unify as switching circuits and digital logic (AND/OR/NOT gates realize connectives, the PDNF gives a sum-of-products implementation, completeness of $\{\uparrow\}$ shows NAND gates alone build every combinational circuit, and SAT solvers serve hardware verification and planning). Relations model relational databases and dependency/grouping (transitive closure answers ``reachable'' or ``manager of a manager''). Groups, rings and fields support coding and cryptography ($\mathbb Z_n$ is the stage of RSA and discrete logarithms, and linear codes are subspaces of $\mathbb F_2^n$). Graph theory supports shortest paths, minimum spanning trees, max-flow, bipartite matching and colouring. For interviews, answer in three layers: language (logic and sets), structure (relations, partial orders, equivalence classes) and algorithms (closures, matching, spanning trees, colouring, topological sorting), giving one concrete term per layer.}

\clearpage
\section*{中英术语表 | Chinese--English glossary}
\noindent\footnotesize
\begin{longtable}{@{}>{\raggedright\arraybackslash}p{3.0cm}>{\raggedright\arraybackslash}p{4.9cm}>{\raggedright\arraybackslash}p{7.0cm}@{}}
\toprule
\textbf{中文} & \textbf{English} & \textbf{简要说明 / Note} \\
\midrule
\endhead
命题 & proposition & 或真或假且不能兼得的陈述句 \\
\hline
重言式 & tautology & 所有赋值下都为真的公式 \\
\hline
主析取范式 & principal disjunctive normal form & 极小项之析取，唯一 \\
\hline
全功能集 & functionally complete set & 能表示一切命题公式的联结词集 \\
\hline
谓词 & predicate & 带空位的断言，如 $P(x)$ \\
\hline
量词 & quantifier & $\forall$ 全称、$\exists$ 存在 \\
\hline
二元关系 & binary relation & 笛卡儿积的子集 \\
\hline
关系的合成 & composition of relations & 先 $R$ 后 $S$ 得到 $R\circ S$ \\
\hline
传递闭包 & transitive closure & 包含 $R$ 的最小传递关系，Warshall 可求 \\
\hline
偏序 & partial order & 自反、反对称、传递的关系 \\
\hline
哈斯图 & Hasse diagram & 只画覆盖关系的偏序图 \\
\hline
等价关系 & equivalence relation & 自反、对称、传递的关系 \\
\hline
等价类 & equivalence class & $[a]=\{x\mid x\sim a\}$ \\
\hline
划分 & partition & 非空、两两不交、并为全集的块族 \\
\hline
代数结构 & algebraic structure & 载体、运算与常元 \\
\hline
同态 & homomorphism & 保持运算的映射 \\
\hline
同余关系 & congruence relation & 与运算相容的等价关系 \\
\hline
群 & group & 封闭、结合、有单位元与逆元 \\
\hline
循环群 & cyclic group & 由单个元素生成的群 \\
\hline
陪集 & coset & $aH=\{ah\mid h\in H\}$ \\
\hline
拉格朗日定理 & Lagrange's theorem & $|G|=|H|\cdot[G:H]$ \\
\hline
域 & field & 每个非零元可逆的含幺交换环 \\
\hline
格 & lattice & 任意二元有上确界与下确界 \\
\hline
分配格 & distributive lattice & 满足分配律的格 \\
\hline
布尔代数 & Boolean algebra & 有补分配格，命题逻辑的代数模型 \\
\hline
二部图 & bipartite graph & 顶点可分成两类，边只跨类 \\
\hline
匹配 & matching & 两两不共端点的边集 \\
\hline
平面图 & planar graph & 可画在平面上使边只在端点相交 \\
\hline
欧拉公式 & Euler's formula & 连通平面图 $V-E+F=2$ \\
\hline
树 & tree & 连通无回路的无向图 \\
\hline
最小生成树 & minimum spanning tree & 总权最小的生成树，Kruskal/Prim 可求 \\
\hline
\end{longtable}
\normalsize
\section*{参考文献 | References}
\small
\noindent 方世昌．离散数学（第三版）．西安：西安电子科技大学出版社，2009．第 1 章（数理逻辑）、第 3 章（二元关系）、第 6 章（代数）、第 7 章（格与布尔代数）、第 8 章（图论）．\par
\noindent 左孝凌，李为鑑，刘永才．离散数学．上海：上海科学技术文献出版社，1982．第 1--3 章（命题逻辑、谓词逻辑、集合与关系）、第 6--8 章（代数结构、格与布尔代数、图论）．\par
\noindent K. H. Rosen. \emph{Discrete Mathematics and Its Applications}, 8th ed. New York: McGraw-Hill, 2019. Chapters 1 (logic and proofs), 9 (relations), 10 (graphs), 11 (trees), 12 (Boolean algebra).\par
\EN{Main source: Fang Shichang, \emph{Discrete Mathematics} (3rd ed.), Xidian University Press, 2009, Chapters 1, 3, 6, 7 and 8. Supplementary: Zuo Xiaoling et al., \emph{Discrete Mathematics}, 1982, Chapters 1--3 and 6--8. English reference: Rosen, \emph{Discrete Mathematics and Its Applications}, 8th ed., Chapters 1, 9--12. Chapter-level references; page numbers omitted.}
\end{document}
